\documentclass[11pt,a4paper]{article}
\usepackage[margin=2.4cm]{geometry}
\usepackage{amsmath,amssymb,amsthm,mathtools}
\usepackage[protrusion=true,expansion=false]{microtype}
\usepackage{booktabs,enumitem}
\usepackage{placeins}
\usepackage{tikz}
\usetikzlibrary{arrows.meta,positioning,calc,fit,backgrounds}
\definecolor{tierspend}{RGB}{200,60,60}
\definecolor{tierfull}{RGB}{210,140,40}
\definecolor{tierin}{RGB}{60,130,180}
\definecolor{tierrecv}{RGB}{70,150,90}
\usepackage[colorlinks=true,linkcolor=blue!50!black,citecolor=blue!50!black,urlcolor=blue!50!black]{hyperref}

\emergencystretch=3em
\hbadness=10000
\hfuzz=2pt

\newtheorem{theorem}{Theorem}[section]
\newtheorem{lemma}[theorem]{Lemma}
\newtheorem{proposition}[theorem]{Proposition}
\newtheorem{corollary}[theorem]{Corollary}
\theoremstyle{definition}
\newtheorem{definition}[theorem]{Definition}
\newtheorem{construction}[theorem]{Construction}
\newtheorem{assumption}[theorem]{Assumption}
\newtheorem{remark}[theorem]{Remark}
\newtheorem{example}[theorem]{Example}

\newcommand{\F}{\mathbb{F}}
\newcommand{\Z}{\mathbb{Z}}
\newcommand{\N}{\mathbb{N}}
\newcommand{\G}{\mathbb{G}}
\newcommand{\E}{\mathbb{E}}
\newcommand{\PP}{\mathbb{P}}
\renewcommand{\Pr}{\mathrm{Pr}}
\newcommand{\Fp}{\F_p}
\newcommand{\Fq}{\F_q}
\newcommand{\ip}[2]{\langle #1,\, #2 \rangle}
\newcommand{\Comm}{\mathsf{Commit}}
\newcommand{\Open}{\mathsf{Open}}
\newcommand{\Setup}{\mathsf{Setup}}
\newcommand{\Verify}{\mathsf{Verify}}
\newcommand{\negl}{\mathsf{negl}}
\newcommand{\poly}{\mathsf{poly}}
\newcommand{\PPT}{\mathsf{PPT}}
\newcommand{\Adv}{\mathcal{A}}
\newcommand{\Prov}{\mathcal{P}}
\newcommand{\Ver}{\mathcal{V}}
\newcommand{\Sim}{\mathcal{S}}
\newcommand{\Ext}{\mathcal{E}}
\newcommand{\Rel}{\mathcal{R}}
\newcommand{\Lang}{\mathcal{L}}
\newcommand{\nfsym}{\mathsf{nf}}
\newcommand{\cmsym}{\mathsf{cm}}
\newcommand{\cmx}{\mathsf{cmx}}
\newcommand{\cvsym}{\mathsf{cv}}
\newcommand{\rk}{\mathsf{rk}}
\newcommand{\Pallas}{\ensuremath{\mathsf{Pallas}}}
\newcommand{\Vesta}{\ensuremath{\mathsf{Vesta}}}
\newcommand{\Ep}{\ensuremath{\mathcal{E}}}
\newcommand{\NoteCommit}{\ensuremath{\mathsf{NoteCommit}}}
\newcommand{\Sinsemilla}{\ensuremath{\mathsf{Sinsemilla}}}
\newcommand{\ShortCommit}{\ensuremath{\mathsf{ShortCommit}}}
\newcommand{\CommitIvk}{\ensuremath{\mathsf{Commit}^{\mathsf{ivk}}}}
\newcommand{\Extract}{\ensuremath{\mathsf{Extract}}}
\newcommand{\Acc}{\ensuremath{\mathsf{Acc}}}
\newcommand{\GroupHash}{\ensuremath{\mathsf{GroupHash}}}
\newcommand{\ask}{\ensuremath{\mathsf{ask}}}
\newcommand{\aksym}{\ensuremath{\mathsf{ak}}}
\newcommand{\nksym}{\ensuremath{\mathsf{nk}}}
\newcommand{\ivk}{\ensuremath{\mathsf{ivk}}}
\newcommand{\fvk}{\ensuremath{\mathsf{fvk}}}
\newcommand{\ovk}{\ensuremath{\mathsf{ovk}}}
\newcommand{\dk}{\ensuremath{\mathsf{dk}}}
\newcommand{\pkd}{\ensuremath{\mathsf{pk_d}}}
\newcommand{\gd}{\ensuremath{\mathsf{g_d}}}
\newcommand{\rcv}{\ensuremath{\mathsf{rcv}}}
\newcommand{\rcm}{\ensuremath{\mathsf{rcm}}}
\newcommand{\rivk}{\ensuremath{\mathsf{r_{ivk}}}}
\newcommand{\rsk}{\ensuremath{\mathsf{rsk}}}

\renewenvironment{abstract}{%
  \small\begin{center}{\bfseries\abstractname}\end{center}\medskip\noindent\ignorespaces
}{\par\medskip}

\title{\textbf{\Huge FROST Guide}\\[6pt]\large Threshold RedPallas: FROST, its re-randomised form, and where Zcash uses it}
\author{m@rek.onl}
\date{}
\begin{document}
\maketitle
\thispagestyle{empty}
\begin{abstract}
Assuming the RedPallas signatures of the \emph{Ironwood Guide}, this volume
of \emph{The Zcash Arboretum} documents threshold signing for Zcash: FROST as
its paper constructs it and as RFC~9591 standardises it, the re-randomised
form that survives RedPallas key re-randomisation, and the surfaces that use
it --- spend authorisation, recoverability custody, and finalizer keys. The
audit boundary is stated plainly: the deployed re-randomised path has not been
audited, and every claim beyond the paper, the RFC, and the draft ZIP carries
its classification.
\end{abstract}
\clearpage
\tableofcontents
\clearpage
%% FROST Guide, section "Scope, sources, and deployment surfaces": the
%% volume opener --- the RedPallas scope restriction, the pinned source
%% ladder, the three-bin rule, and the three custody surfaces.
%% Body-only LaTeX; labels sec:fr-context, sec:fr-scope, sec:fr-sources,
%% sec:fr-bins, sec:fr-surfaces, tab:fr-sources. Register per
%% ironwood-guide; 80 cols.
%%
%% Sources (all verified firsthand 2026-07-15):
%%   - RFC 9591, "The Flexible Round-Optimized Schnorr Threshold (FROST)
%%     Protocol for Two-Round Schnorr Signatures", Connolly/Komlo/
%%     Goldberg/Wood, IRTF (CFRG stream), Category Informational, ISSN
%%     2070-1721, June 2024. https://www.rfc-editor.org/rfc/rfc9591.html;
%%     local copy .wip/frost-drafts/sources/rfc9591.txt.
%%   - C. P. L. Gouvea, C. Komlo, "Re-Randomized FROST", IACR ePrint
%%     2024/436, https://eprint.iacr.org/2024/436 (received 2024-03-13,
%%     approved 2024-03-15; publication info "Preprint", no venue). Read
%%     in full; local copies sources/rerand-2024-436.pdf/.txt. eprint's
%%     .pdf endpoint intermittently behind a Cloudflare challenge as of
%%     2026-07-15 (abstract page reachable; direct PDF curl 403'd).
%%   - zcash/zips @ 69610984109078075e7e64989ecf89d63a259b97 (2026-07-09):
%%     zip-0312.rst (Draft), zip-2005.md (Proposed), zip-0271.md,
%%     zip-0326.md, zip-0374.md; zip-0313.rst checked and excluded.
%%   - ZcashFoundation/frost @ 2016e44ba4a4757a996300350063b937a2ad33e8
%%     (2026-04-23), workspace 3.0.0, rust-version 1.81.
%%   - ZcashFoundation/reddsa @ 3792daa95e588c1af6bd4805105bfb6ea7e9ad49
%%     (2026-05-22), crate 0.5.2.
%%   - Antecedents: IACR ePrint 2020/852 (Komlo--Goldberg, SAC 2020) and
%%     2022/833 (Bellare--Crites--Komlo--Maller--Tessaro--Zhu, CRYPTO
%%     2022); local copies under sources/.
%%   - Sibling volumes: ironwood-guide.tex sec:redpallas (line 1308),
%%     sec:redpallas-sec (line 1742); wallet-guide.tex sec:pczt (line
%%     2882); crosslink-guide.tex sec:xl-stake-keys (line 2889).
%%
%% Notes:
%%   - The volume brief said "likely zip-0312/0313"; ZIP 313 is the
%%     obsolete fee ZIP (obsoleted by ZIP 317). The FROST ZIP is 312
%%     alone; the second FROST-bearing ZIP is 2005. Footnoted in body.
%%   - The randomizer wire-object discrepancy (paper sends alpha-hat,
%%     ZIP 312 sends the HR output, frost-rerandomized 3.0.0 sends a
%%     seed) is NOT touched here; it is resolved in sec:fr-deploy-rand.

\section{Scope, sources, and deployment surfaces}\label{sec:fr-context}

This volume documents threshold custody of Zcash spend authority: how $t$
of $n$ parties, no one of whom ever holds the whole signing key, jointly
produce a signature that consensus accepts as an ordinary RedPallas
spend-authorisation signature. Three artefacts carry the construction.
The base scheme is FROST, the two-round threshold Schnorr protocol of
Komlo and Goldberg, standardised as RFC 9591 (\S\ref{sec:fr-frost}). The
Zcash layer is Re-Randomized FROST (Gouv\^ea--Komlo, IACR ePrint
2024/436), which shifts every share by a common randomizer so that the
aggregate signature verifies under the per-Action randomised key; ZIP 312
specifies it for Zcash, and the \texttt{frost-rerandomized} and
\texttt{reddsa} crates ship it (\S\ref{sec:fr-rerand}). The deployment
picture --- the running stack, the proposed uses, and the gaps --- closes
the volume (\S\ref{sec:fr-deploy}).

Three custody surfaces anchor the treatment, one per rigor bin
(\S\ref{sec:fr-bins}): spend-authorisation custody for Orchard, where the
construction is specified and shipped; custody of the ZIP-2005 quantum
spending key $\mathsf{qsk}$, where a Proposed consensus ZIP constrains a
FROST ceremony whose recovery protocol it defers; and Crosslink finalizer
keys, where no threshold design exists at all. The surfaces and their
bins are fixed in \S\ref{sec:fr-surfaces}.

\subsection{Scope}\label{sec:fr-scope}

The object under study is a $t$-of-$n$ signing protocol whose output a
verifier cannot tell apart from single-party spend authorisation. ZIP 312
states the requirements directly: the signatures must verify as Sapling or
Orchard spend-authorisation signatures, and should meet the protocol
specification's ``Signature with Re-Randomizable Keys'' criteria. Plain
FROST fails the first clause --- it signs for a fixed group key, and
consensus verifies only per-Action randomised keys --- and closing that
gap is the volume's central construction (\S\ref{sec:fr-rerand}).

The volume develops RedPallas only, per its title. ZIP 312 itself defines
two ciphersuites --- FROST(Jubjub, BLAKE2b-512) for Sapling RedJubjub and
FROST(Pallas, BLAKE2b-512) for Orchard RedPallas --- and we cite the
Jubjub suite where the ZIP fixes both at once, but develop only the
Pallas suite (\S\ref{sec:fr-rerand-suite}).

ZIP 312's stated non-requirements bound the scope from below: it does not
remove the Coordinator role; it defines no key-generation procedure (out
of scope there, exactly as in RFC 9591); and it provides no network
privacy. Key generation nevertheless appears twice in this volume, both
times forced by sources: once as published with the base scheme
(\S\ref{sec:fr-keygen}), and once as the unspecified-but-shipped DKG
surface of the \texttt{reddsa} crate, binned in \S\ref{sec:fr-bins}.

One altitude note governs every deployment claim. The deployment is
library-level, never consensus-level: no consensus rule knows FROST
exists, and none needed to change. The verifier accepts the aggregate
because it is a valid RedPallas signature; as ZIP 312's rationale puts
it, the method by which the commitment $R$ was generated is irrelevant to
the verifier.

\subsection{The source ladder}\label{sec:fr-sources}

Table~\ref{tab:fr-sources} pins every primary source. All repository
facts in this volume were verified firsthand at the stated commits; web
sources were fetched 2026-07-15; local copies of the papers and the RFC
are archived with the volume's drafts.

\begin{table}[htbp]
\centering
\small
\begin{tabular}{@{}lll@{}}
\toprule
Artefact & Pin / date & Rigor \\
\midrule
RFC 9591 (IRTF CFRG)
  & June 2024                      & Informational; standards rigor \\
ePrint 2024/436 (Gouv\^ea--Komlo)
  & recv.\ 2024-03-13              & preprint; no venue \\
ZIP 312
  & \texttt{6961098410}, 2026-07-09 & Draft (Wallet category) \\
ZIP 2005
  & \texttt{6961098410}, 2026-07-09 & Proposed (Consensus category) \\
\texttt{ZcashFoundation/frost}
  & \texttt{2016e44ba4}, 2026-04-23 & shipped code; workspace 3.0.0 \\
\texttt{ZcashFoundation/reddsa}
  & \texttt{3792daa95e}, 2026-05-22 & shipped code; crate 0.5.2 \\
ePrint 2020/852; 2022/833
  & SAC 2020; CRYPTO 2022          & peer-reviewed antecedents \\
\bottomrule
\end{tabular}
\caption{Primary FROST sources, pinned. Verification date 2026-07-15.}
\label{tab:fr-sources}
\end{table}

\paragraph{The standard.} RFC 9591, ``The Flexible Round-Optimized
Schnorr Threshold (FROST) Protocol for Two-Round Schnorr Signatures''
(D. Connolly, Zcash Foundation; C. Komlo, University of Waterloo and
Zcash Foundation; I. Goldberg, University of Waterloo; C. A. Wood,
Cloudflare; IRTF CFRG stream, Category Informational, June 2024), is the
top of the ladder and the only artefact here at standards rigor. Two of
its four authors are Zcash Foundation: the standard and the Zcash
deployment share provenance. It fixes the base two-round protocol and
leaves key generation out of scope; its deltas against the SAC 2020 paper
are catalogued in \S\ref{sec:fr-rfc}.

\paragraph{The paper.} ``Re-Randomized FROST'', by Gouv\^ea and Komlo
(both Zcash Foundation), is IACR ePrint 2024/436
(\texttt{https://eprint.iacr.org/2024/436}): received 2024-03-13,
approved 2024-03-15, publication info ``Preprint'', no peer-reviewed
venue as of 2026-07-15. We read it in full and work from it throughout
\S\ref{sec:fr-rerand}. Its headline result, Theorem~1, proves
unforgeability in the random-oracle model under AOMDL --- the same
assumptions as plain FROST (\S\ref{sec:fr-rerand-sec}). The construction
it describes is specified and shipped; the security theorem above it
rests on a preprint, and we say so wherever we lean on it.

\paragraph{The ZIP.} ZIP 312, ``FROST for Spend Authorization
Multisignatures'' (owners Conrado Gouv\^ea, Chelsea Komlo, and Deirdre
Connolly; Status Draft, Category Wallet; Discussions-To zcash/zips\#382,
pull request zcash/zips\#662), is the Zcash specification of the
construction.\footnote{The FROST ZIP is ZIP 312 alone. ZIP 313, its
numeric neighbor, is ``Reduce Conventional Transaction Fee to 1000
zatoshis'' --- Status Obsolete, obsoleted by ZIP 317 --- and has nothing
to do with FROST.} It is visibly unfinished --- its Created field still
reads \texttt{2022-08-dd}, day unfilled --- so we cite it as a moving
object, by repository commit (\texttt{zcash/zips} @
\texttt{6961098410}, 2026-07-09) and Draft status, never as a stable
specification. Its threat model trusts the Coordinator and every
key-share holder with transaction privacy and unlinkability, never with
security: a rogue Coordinator cannot create a signed transaction without
the approval of \texttt{MIN\_PARTICIPANTS} participants. What it
specifies --- the unlinkability requirement, the randomizer flow, and the
two ciphersuites --- is developed in \S\ref{sec:fr-rerand-zip} and
\S\ref{sec:fr-rerand-suite}.

\paragraph{The crates.} Repository \texttt{ZcashFoundation/frost}
(@ \texttt{2016e44ba4}, 2026-04-23) is a Rust workspace at version 3.0.0
(\texttt{rust-version} 1.81): crate \texttt{frost-core} carries the
ciphersuite-generic protocol, six ciphersuite crates instantiate it
(ed25519, ed448, P-256, ristretto255, secp256k1, secp256k1-tr),
\texttt{frost-rerandomized} carries the re-randomised layer, and
\texttt{gencode} generates the per-suite code. The repository also
vendors an audit report dated 2021-03-23
(\texttt{zcash-frost-audit-report-20210323.pdf}) and the source of the ZF
FROST book (published at \texttt{frost.zfnd.org}), whose ``FROST with
Zcash'' chapters (zcash-devtool demo, Ywallet demo, FROST server) and
serialisation-format page --- the page ZIP 312 cites for
\texttt{encode\_signing\_package} --- feed \S\ref{sec:fr-deploy}.
Repository \texttt{ZcashFoundation/reddsa} (@ \texttt{3792daa95e},
2026-05-22; crate version 0.5.2) holds the Zcash ciphersuites: its
\texttt{frost} feature enables \texttt{frost-rerandomized} 3.0.0, and
\texttt{src/frost/redpallas.rs} defines FROST(Pallas, BLAKE2b-512) as
struct \texttt{PallasBlake2b512}, ZIP 312's reference implementation
surface (\S\ref{sec:fr-rerand-impl}). Crate \texttt{reddsa} also ships
machinery ZIP 312 does not specify --- DKG hash personalisations and
even-$Y$ normalisation hooks --- binned in \S\ref{sec:fr-bins}.

\paragraph{The antecedent papers.} FROST itself is Komlo--Goldberg, IACR
ePrint 2020/852 (SAC 2020); the security analysis the re-randomised
extraction argument extends is Bellare--Crites--Komlo--Maller--Tessaro--%
Zhu, ``Better than Advertised Security for Non-Interactive Threshold
Signatures'' (CRYPTO 2022; ePrint 2022/833). They are the paper's
references [10] and [1] respectively; \S\ref{sec:fr-frost} works from
both.

\subsection{The three-bin rule}\label{sec:fr-bins}

The series rule is binary for deployed behaviour and ternary for design.
A claim about deployed behaviour cites the implementation --- crate and
file --- and the specification section it realises; a claim about a
design-stage protocol is classified \emph{specified}, \emph{designed but
unspecified}, or an \emph{open problem}, and the bin stays visible in the
prose. This volume sits astride the boundary: one surface is shipped
code, one is a consensus proposal, one is an acknowledged gap. Applied
once, here, for the whole volume:
\begin{enumerate}
\item Re-Randomized FROST for spend-authorisation custody is
  \emph{specified}: RFC 9591 published June 2024, ZIP 312 in Draft,
  ePrint 2024/436, and shipped code --- \texttt{frost-rerandomized}
  3.0.0 and \texttt{reddsa} 0.5.2. The deployment is library-level, not
  consensus (\S\ref{sec:fr-scope}).
\item Custody of the ZIP-2005 quantum spending key is specified at the
  ZIP level (Status Proposed, Category Consensus), but the Recovery
  Protocol it exists to serve is ``to be introduced, potentially, in a
  future upgrade'' --- \emph{designed but unspecified}.
\item Crosslink finalizer threshold custody is an \emph{open problem}:
  the \emph{Crosslink Guide}'s own classification, which we adopt.
\end{enumerate}

One surface refuses the clean split, and we bin it here once. No ZIP
specifies FROST key generation: ZIP 312 declares it out of scope (as does
RFC 9591), and ZIP 271's May-2025 note confirms the gap --- ZIP 312's
missing key-generation text is its stated reason not to couple the
lockbox disbursement to that specification. Yet crate \texttt{reddsa}
ships DKG support: hash personalisations \texttt{FROST\_RedPallasD}
(HDKG) and \texttt{FROST\_RedPallasI} (HID), beyond ZIP 312's H1--H5 and
HR, plus even-$Y$ normalisation hooks
(\texttt{post\_generate}/\texttt{post\_dkg}, applying
\texttt{into\_even\_y} to secret shares and the public key package) so
that the group verifying key meets the encoding constraint of Orchard's
spend validating key. Deployed code with no specification above it: we
classify the DKG surface \emph{designed but unspecified} and cite the
crate at commit \texttt{3792daa95e} wherever we rely on it.
%% DKG bin confirmed 2026-07-15: zips PR #895 (key generation for ZIP
%% 312) is still open (created 2024-08-13, 21 commits, awaiting review
%% by daira and nuttycom), so no ZIP or RFC specifies any DKG for Zcash
%% custody; #895's merge is the revisit trigger.

\subsection{The three deployment surfaces}\label{sec:fr-surfaces}

\paragraph{Spend-authorisation custody (specified).} The product ZIP 312
promises is signatures ``indistinguishable from RedPallas Orchard Spend
Authorization Signatures''. The single-party object being imitated is
settled in the \emph{Ironwood Guide}: \S``Authorisation of the spend:
RedPallas'' constructs $\rsk := \ask + \alpha$ and
$\rk := \aksym + [\alpha]\,G^{\mathsf{SpendAuth}} =
[\rsk]\,G^{\mathsf{SpendAuth}}$, with $\mathsf{GenRandom}$
hashing $80$ uniform bytes through $H^{\circledast}$, and \S``RedPallas
unforgeability and re-randomisation'' proves EUF-CMA from discrete log in
the random-oracle model. Per the layering rule we cite both and re-prove
neither; this volume's work begins where the single signer ends, at
splitting $\ask$ into shares no signing coalition ever
reassembles (\S\ref{sec:fr-rerand}). The wallet-side integration path is
likewise on the record: the \emph{Wallet Guide}, \S``Partially created
transactions (PCZT)'', records ZIP 374's motivation (iii), ``FROST
threshold signing (RFC 9591)'' --- the multiparty transaction format
FROST signing rides in \S\ref{sec:fr-deploy-pczt}.

\paragraph{Custody of the quantum spending key (designed but
unspecified).} ZIP 2005, ``Ironwood Quantum Recoverability'' (owners
Daira-Emma Hopwood and Jack Grigg; Status Proposed, Category Consensus;
created 2025-03-31), writes FROST into its requirements twice: the design
should be fully compatible with FROST threshold multisignatures (citing
ZIP 312), and should lose no pre-quantum security when FROST and/or
hardware wallets are used. Its ``Usage with FROST'' section then
constrains the ceremony: the key $\aksym$ is derived jointly, via
the FROST DKG, from shares of $\ask$; the participants must
privately agree on a value $\mathsf{sk}$ and use it with
$\mathsf{use\_qsk} = \mathsf{true}$ to derive $\nksym$,
$\mathsf{qsk}$, $\mathsf{qk}$, and $\mathsf{rivk\_ext} =
\mathsf{H}^{\mathsf{rivk\_ext}}_{\mathsf{qk}}(\aksym, \nksym)$;
all participants must obtain the same $\aksym$, $\nksym$, and
$\mathsf{rivk\_ext}$; and each participant must hold $\mathsf{qsk}$ at
least as securely as $\ask$. The Recovery Protocol requires a
RedDSA signature under $\aksym$ --- preserving $t$-of-$n$ for as
long as discrete log on Pallas stands --- together with a proof
$\mathrm{SoK}^{\mathsf{qsk}}$ of knowledge of a $\mathsf{qsk}$ with
$\mathsf{H}^{\mathsf{qk}}(\mathsf{qsk}) = \mathsf{qk}$. The asymmetry is
the point to retain: the key $\mathsf{qsk}$ is held by \emph{every}
participant, so a quantum adversary needs only $\mathsf{qsk}$ and the
threshold property is gone; the ZIP accordingly recommends moving funds
to a fully post-quantum threshold protocol once one is available. The
deferred Recovery Protocol is what parks this surface in the middle bin;
the full treatment is \S\ref{sec:fr-deploy-qsk}.

\paragraph{Crosslink finalizer keys (open problem).} The \emph{Crosslink
Guide}, \S``Finalizer keys'', records the prototype state: one ed25519
key per finalizer (\texttt{ed25519-zebra}), a $76$-byte vote format
signed by appending a $64$-byte ed25519 signature over those $76$
bytes, and --- by negative search over the four
Crosslink repositories (tfl-book, crosslink-protocol-guide,
crosslink-spec, zebra-crosslink), dated 2026-07-15 there --- no FROST or
threshold-signing design for finalizer keys sourced anywhere in the
corpus. That guide's classification bins threshold custody, key
rotation, revocation, and initial-roster formation as open problems, and
observes that FROST's presence in the PCZT flow makes the absence ``a
visible gap rather than an exotic ask''. We adopt the classification and
return to the surface in \S\ref{sec:fr-deploy-crosslink}.

\paragraph{Adjacent surfaces, cited but not developed.} Three further
ZIPs touch FROST and mark how wide the surface is growing. ZIP 271
(``Dev Fund Extension and One-Time Disbursement'', Proposed) anticipates
spending lockbox one-time-disbursement chunks ``to a FROST address'',
and its May-2025 note that ZIP 312 does not specify key generation is
its stated reason not to couple the two specifications. ZIP 326
(``NU6.3 Consequences for Wallets'', Draft, created 2026-06-30) cites
ZIP 2005's ``Usage with FROST'' section. ZIP 374 (``Partially Created
Zcash Transaction Format'', Draft, created 2024-12-09) lists FROST per
RFC 9591 under its multiparty-transactions motivation. Each appears in
this volume exactly where it bears (\S\ref{sec:fr-deploy}); none is
developed as a surface of its own.
%% TODO(cross-refs): once this volume lands, the Wallet Guide's sec:pczt
%% and the Crosslink Guide's sec:xl-stake-keys prose --- both already
%% citing FROST/RFC 9591 --- become candidates for italic-name citations
%% of the FROST Guide; schedule in the adversarial review's
%% cross-reference check.

%% FROST Guide, section "FROST": the scheme as published (Komlo--Goldberg,
%% ePrint 2020/852), the corrected later analyses, and the RFC 9591 deltas.
%% Body-only LaTeX; labels sec:fr-*. Register per ironwood-guide.
%%
%% Sources (all verified firsthand 2026-07-15):
%%   - C. Komlo, I. Goldberg, "FROST: Flexible Round-Optimized Schnorr
%%     Threshold Signatures", IACR ePrint 2020/852,
%%     https://eprint.iacr.org/2020/852 (received 2020-07-12, last revised
%%     2020-12-22 -- the version read, in full; published at SAC 2020,
%%     LNCS 12804, pp. 34-65). PDF fetched via Wayback (eprint serves a
%%     Cloudflare challenge to curl): .wip/frost-drafts/sources/
%%     frost-2020-852.pdf/.txt.
%%   - M. Bellare, S. Tessaro, C. Zhu, "Stronger Security for
%%     Non-Interactive Threshold Signatures: BLS and FROST", IACR ePrint
%%     2022/833 (June 2022; preliminary version in Bellare-Crites-Komlo-
%%     Maller-Tessaro-Zhu, "Better than Advertised Security for
%%     Non-interactive Threshold Signatures", CRYPTO 2022).
%%   - C. P. L. Gouvea, C. Komlo, "Re-Randomized FROST", IACR ePrint
%%     2024/436, https://eprint.iacr.org/2024/436 (received 2024-03-13).
%%     Theorem 1 bound checked against the PDF page render (p9-09.png,
%%     paper p. 9, Equation 3): the radical spans all three terms,
%%     including "- negl(lambda)".
%%   - RFC 9591, "The Flexible Round-Optimized Schnorr Threshold (FROST)
%%     Protocol for Two-Round Schnorr Signatures", Connolly/Komlo/
%%     Goldberg/Wood, IRTF CFRG, June 2024.
%%   - ZIP 312 (Draft), ~/zcash/zips/zips/zip-0312.rst, repo commit
%%     6961098410 (2026-07-09).
%%   - ZcashFoundation/frost, commit 2016e44ba4 (2026-04-23), workspace
%%     version 3.0.0; ZcashFoundation/reddsa, commit 3792daa95e
%%     (2026-05-22), version 0.5.2.
%%
%% Notation note: this section retains the source paper's multiplicative
%% notation (g^x, capital-letter group elements) so that every displayed
%% equation is verbatim from ePrint 2020/852; the bridge to the series'
%% additive notation ([k]P on Pallas) is made once, in the RedPallas
%% instantiation section. Policy confirmed at the 2026-07-15 review:
%% every displayed equation in sec:fr-frost and def:fr-rerand-frost was
%% checked character-for-character against the sources, which is what
%% the source-verbatim multiplicative policy exists to buy; quoted
%% material stays verbatim, and all volume-level constructions use the
%% series-additive notation.

\section{FROST}\label{sec:fr-frost}

The base object of this volume is the threshold Schnorr scheme FROST
(\emph{Flexible Round-Optimized Schnorr Threshold} signatures) of Komlo and
Goldberg, published at SAC 2020 and maintained as IACR ePrint 2020/852; we
work from the final revision of 22 December 2020. The scheme lets any $t$
of $n$ parties, each holding a Shamir share of a signing key $s$, produce a
signature $(R, z)$ that verifies under the ordinary single-party Schnorr
equation for the group public key $Y = g^s$ --- a verifier cannot
distinguish a FROST signature from a single-signer one. That
indistinguishability is precisely what makes the scheme deployable under
Zcash's RedDSA validators: the consensus rules never learn that a threshold
ceremony occurred. The deployment surfaces this volume targets are
spend-authorisation custody over RedPallas (the \emph{Ironwood Guide},
\S``Authorisation of the spend: RedPallas''), custody of the ZIP-2005
quantum spending key (whose ``Usage with FROST'' section constrains the key
generation), and Crosslink finalizer keys (the \emph{Crosslink Guide},
\S``Finalizer keys'').

Throughout, $\G$ is a group of prime order $q$ with generator $g$, and
$H$, $H_1$, $H_2$ are hash functions with outputs in $\Z_q^*$, modelled as
random oracles. The paper fixes the single-party scheme it must match
--- the Schnorr signature scheme of the \emph{Crypto Guide}, \S``From
identification to signature via the Fiat--Shamir transform'' ---
restated here in the paper's multiplicative notation: a Schnorr
signature over $m$ under secret key $s \in \Z_q$ and
public key $Y = g^s \in \G$ is produced by sampling a nonce
$k \xleftarrow{\$} \Z_q$, setting $R = g^k$, $c = H(R, Y, m)$,
$z = k + s \cdot c$, and outputting $\sigma = (R, z)$; verification
recomputes $c$ and accepts iff $R = g^z \cdot Y^{-c}$. FROST's whole design
problem is to compute the pair $(R, z)$ jointly, with $k$ and $s$ existing
only as shares, without opening the concurrency attacks described in
\S\ref{sec:fr-binding}.

\subsection{Shamir sharing and share conversion}\label{sec:fr-shamir}

Two secret-sharing mechanisms cooperate in FROST, and keeping them distinct
is the key to reading the signing equation.

\paragraph{Shamir sharing (the long-lived key).} A dealer --- or, in
\S\ref{sec:fr-keygen}, every participant acting as a dealer simultaneously
--- selects $t-1$ random coefficients $a_1, \dots, a_{t-1}$ and forms the
degree-$(t-1)$ polynomial
\[
f(x) = s + \sum_{i=1}^{t-1} a_i x^i, \qquad f(0) = s,
\]
giving participant $P_i$ the share $(i, f(i))$. Any $t$ shares determine
$f$ by Lagrange interpolation and hence $s = f(0)$; fewer than $t$ reveal
nothing. For a signing set $S = \{p_1, \dots, p_\alpha\}$ of $\alpha$
participant identifiers ($t \le \alpha \le n$), the interpolation weight of
participant $i$ at zero is the Lagrange coefficient
\[
\lambda_i = \prod_{j=1, j \ne i}^{\alpha} \frac{p_j}{p_j - p_i},
\]
computable from the identifiers alone --- the coalition can be chosen
\emph{after} shares are dealt, which is what makes FROST's ``any $t$ of
$n$, decided at signing time'' flexibility possible.

\paragraph{Additive sharing and conversion (the per-signature nonce).} A
sum $s = \sum_{i} s_i$ with each summand chosen locally is an additive
sharing: it needs no dealer and no interaction. The
Cramer--Damg{\aa}rd--Ishai share conversion observation, used by FROST in
its simplest case, is that additive shares convert locally to Shamir form:
if $s = \sum_{i \in S} s_i$, then the values $s_i / \lambda_i$ are Shamir
shares of the same $s$ (evaluations of a degree-$(\alpha-1)$ polynomial
interpolating them). FROST uses exactly this to make the group nonce $k$:
each signer contributes an additively-shared summand non-interactively,
and the $\lambda_i$ weights in the signing equation perform the conversion
so that the response aggregates by plain summation. No DKG runs at signing
time.

\subsection{Key generation: Pedersen's DKG with a knowledge proof}
\label{sec:fr-keygen}

FROST generates its long-lived shares with a variant of Pedersen's DKG ---
$n$ parallel executions of Feldman's verifiable secret sharing, one with
each participant as dealer --- hardened with a proof of knowledge of each
dealer's free coefficient. The paper's Figure 1, reproduced:

\paragraph{Round 1.}
\begin{enumerate}
\item Every participant $P_i$ samples $t$ random values
  $(a_{i0}, \dots, a_{i(t-1)}) \xleftarrow{\$} \Z_q$ and defines the
  degree-$(t-1)$ polynomial $f_i(x) = \sum_{j=0}^{t-1} a_{ij} x^j$.
\item Every $P_i$ computes a proof of knowledge of $a_{i0}$: with
  $k \xleftarrow{\$} \Z_q$, $R_i = g^k$,
  $c_i = H(i, \Phi, g^{a_{i0}}, R_i)$,
  $\mu_i = k + a_{i0} \cdot c_i$, set $\sigma_i = (R_i, \mu_i)$, where
  $\Phi$ is a context string preventing replay across ceremonies.
\item Every $P_i$ computes the public commitment
  $\vec{C}_i = \langle \phi_{i0}, \dots, \phi_{i(t-1)} \rangle$ with
  $\phi_{ij} = g^{a_{ij}}$.
\item Every $P_i$ broadcasts $\vec{C}_i, \sigma_i$.
\item On receiving $\vec{C}_\ell, \sigma_\ell$, participant $P_i$ verifies
  $R_\ell \stackrel{?}{=} g^{\mu_\ell} \cdot \phi_{\ell 0}^{-c_\ell}$ with
  $c_\ell = H(\ell, \Phi, \phi_{\ell 0}, R_\ell)$, \emph{aborting on
  failure}; on success all $\sigma_\ell$ are deleted.
\end{enumerate}

\paragraph{Round 2.}
\begin{enumerate}
\item Each $P_i$ securely sends each other participant $P_\ell$ the share
  $(\ell, f_i(\ell))$, then deletes $f_i$ and every share except its own
  $(i, f_i(i))$.
\item Each $P_i$ verifies each received share against the dealer's
  commitment:
  \[
  g^{f_\ell(i)} \stackrel{?}{=} \prod_{k=0}^{t-1} \phi_{\ell k}^{\,i^k
  \bmod q},
  \]
  aborting if the check fails.
\item Each $P_i$ computes its long-lived signing share
  $s_i = \sum_{\ell=1}^{n} f_\ell(i)$, stores it securely, and deletes the
  $f_\ell(i)$.
\item Each $P_i$ computes its public verification share $Y_i = g^{s_i}$
  and the group public key $Y = \prod_{j=1}^{n} \phi_{j0}$.
\end{enumerate}

Three design decisions deserve comment. First, the \emph{proof of
knowledge of $a_{i0}$} is the delta over plain Pedersen DKG: it forecloses
rogue-key attacks in the $t \ge n/2$ regime, where a participant could
otherwise choose its commitment as a function of others' to bias the
aggregate key (the fix was suggested to the authors by Shlomovits and
Turkel; the paper's changelog dates it to the 2020-07-08 revision).
Second, the \emph{complaint handling} is abort-only: where Gennaro et
al.\ add a complaint/disqualification phase to make the DKG robust, FROST
aborts on any failed share check and expects the cause to be investigated
out of band; the paper notes that implementations wanting robustness can
substitute the Gennaro et al.\ construction. Third, a \emph{consistent
view} of the broadcast commitments $\vec{C}_i$ across participants is
assumed, not provided --- the paper defers to whatever broadcast mechanism
the deployment supplies. Each participant ends holding $(i, s_i)$, with
$Y_i = g^{s_i}$ public for share verification during signing and $Y$ the
single key the outside world sees.

\subsection{Preprocessing and the two-round signing protocol}
\label{sec:fr-sign}

Signing splits into a message-independent \emph{preprocess} stage and a
single online round (equivalently, run as one two-round protocol with the
commitment exchange inlined). A semi-trusted \emph{signature aggregator}
$\mathrm{SA}$ coordinates: it is trusted only for liveness and correct
misbehaviour reporting, not for unforgeability --- a malicious
$\mathrm{SA}$ can at worst deny service or falsely accuse, never learn the
private key or cause improper messages to be signed.

\paragraph{Preprocess($\pi$) (paper Figure 2).} Each participant $P_i$
generates $\pi$ single-use nonce pairs and publishes their commitments:
for $1 \le j \le \pi$,
\[
(d_{ij}, e_{ij}) \xleftarrow{\$} \Z_q^* \times \Z_q^*, \qquad
(D_{ij}, E_{ij}) = (g^{d_{ij}}, g^{e_{ij}}),
\]
appending $(D_{ij}, E_{ij})$ to a list $L_i$ published to a predetermined
location, and storing the secret halves for later.

\paragraph{Sign($m$) (paper Figure 3).} The aggregator $\mathrm{SA}$
selects a signing set $S$ of $\alpha$ participants
($t \le \alpha \le n$), fetches each member's next unused commitment
pair, and forms the ordered list
$B = \langle (i, D_i, E_i) \rangle_{i \in S}$, sending $(m, B)$ to every
$P_i, i \in S$. Each participant then:
\begin{enumerate}
\item validates $m$ (a signer signs only messages it is willing to sign)
  and checks $D_\ell, E_\ell \in \G^*$ for every commitment in $B$,
  aborting on failure;
\item computes the \emph{binding values}
  \[
  \rho_\ell = H_1(\ell, m, B), \quad \ell \in S,
  \]
  the group commitment and challenge
  \[
  R = \prod_{\ell \in S} D_\ell \cdot (E_\ell)^{\rho_\ell}, \qquad
  c = H_2(R, Y, m);
  \]
\item computes its response, converting its Shamir share to the additive
  aggregation domain via $\lambda_i$ (over the set $S$):
  \[
  z_i = d_i + (e_i \cdot \rho_i) + \lambda_i \cdot s_i \cdot c ;
  \]
\item securely deletes $((d_i, D_i), (e_i, E_i))$, then returns $z_i$ to
  $\mathrm{SA}$.
\end{enumerate}
Finally $\mathrm{SA}$ recomputes $\rho_i$, $R_i = D_i \cdot
(E_i)^{\rho_i}$, $R = \prod_{i \in S} R_i$ and $c$, and verifies every
share against the public data:
\[
g^{z_i} \stackrel{?}{=} R_i \cdot Y_i^{\,c \cdot \lambda_i},
\]
identifying and reporting the misbehaving participant and aborting if any
check fails; otherwise it publishes $\sigma = (R, z)$ with
$z = \sum_{i \in S} z_i$.

\paragraph{Correctness.} The proof is two interpolations (paper \S6.1).
The key polynomial $F_1(x) = \sum_{j=1}^n f_j(x)$ from key generation has
$s_i = F_1(i)$ and $s = F_1(0)$. The nonce values
$(i, (d_i + e_i \cdot \rho_i)/\lambda_i)$ define a degree-$(\alpha-1)$
polynomial $F_2(x)$ with
$F_2(0) = \sum_{i \in S} d_i + e_i \cdot \rho_i = k$. Setting
$F_3(x) = F_2(x) + c \cdot F_1(x)$, each response is
$z_i = \lambda_i F_3(i)$, so $z = \sum_{i \in S} z_i$ interpolates
$F_3(0) = k + c \cdot s$; with $R = g^k$ and $c = H_2(R, Y, m)$, the pair
$(R, z)$ is a correct Schnorr signature on $m$.

\subsection{Why the binding factor exists: Drijvers and ROS}
\label{sec:fr-binding}

The one non-obvious ingredient is $\rho_i = H_1(i, m, B)$. The binding
factor is FROST's answer to a pair of concurrency attacks that broke or
bounded every prior two-round scheme.

\paragraph{The Drijvers attack.} Against two-round Schnorr multisignatures
in which the adversary may open $T$ parallel sessions, Drijvers et al.\
showed how to forge by finding a hash output that is a sum of others:
$c^* = H(R^*, Y, m^*)$ with $c^* = \sum_{j=1}^{T} H(R_j, Y, m_j)$, using
Wagner's $k$-tree algorithm for the generalised birthday problem in time
$O(\kappa \cdot b \cdot 2^{b/(1 + \lg \kappa)})$ for $\kappa = T + 1$ and
$b$ the bitlength of the group order. The attack transfers from the
$n$-of-$n$ multisignature setting to a $t$-of-$n$ threshold scheme with an
adversary controlling up to $t-1$ signers, and it is why the Gennaro et
al.\ scheme must cap its signing parallelism. Benhamouda et al.'s
subsequent polynomial-time ROS solver (for $\ell > \lg q$ parallel
sessions) turned the same structural weakness from subexponential to
practical against schemes including two-round MuSig and GJKR-DKG-based
threshold Schnorr; their paper concludes that ``the current version of''
FROST is \emph{not} affected.

\paragraph{How binding forecloses it.} Because each signer derives
$\rho_i = H_1(i, m, B)$ before responding, its share $z_i$ is bound to
this message, this signing set, and this exact commitment list. Responses
cannot be recombined across sessions, messages, or reordered signer sets:
any such mix changes some $\rho_\ell$, hence $R$, hence $c$, and the
transplanted share fails verification. The paper's own generalisation (its
Equation 1, with linear combinations $\gamma_j$) states the residual
attack surface precisely: the forger must find
\[
H_2(R^*, Y, m^*) = \sum_{j} \gamma_j \cdot H_2(R_j, Y, m_j),
\]
where $R^* = g^{r^*} \cdot \prod_j \widehat{R}_j^{\gamma_j}$ is itself a
function of every input on the right-hand side. Wagner's algorithm needs
the two sides of the collision to vary independently; here every candidate
combination on the right moves the left-hand target, degrading the
generalised birthday collision into multiple preimage attacks. The paper's
footnote 7 concedes this step is heuristic: sufficiency of the condition
is immediate, necessity is not proved.

\subsection{Nonce hygiene}\label{sec:fr-nonces}

The preprocessed pairs $(d_{ij}, e_{ij})$ are one-time keys in the
strictest sense. The paper's stipulations:
\begin{itemize}
\item each pair is used at most once, and the signer \emph{securely
  deletes} it in the same signing step that returns $z_i$ (Figure 3,
  step 6);
\item an accidentally reused $(d_{ij}, e_{ij})$ can expose the long-term
  share $s_i$;
\item a $\mathrm{SA}$ that maliciously replays a commitment pair achieves
  nothing: the participant simply aborts, so replay grants no power;
\item the commitment store (if a server) is trusted only for availability
  and freshness; possession of the public $(D_{ij}, E_{ij})$ lists grants
  no signing power.
\end{itemize}
See \S\ref{sec:fr-rfc} for RFC 9591's sharpening of nonce generation
(hedged sampling, and an explicit warning that deterministic nonces enable
complete key recovery in the multi-party setting).

\subsection{The security theorem as stated}\label{sec:fr-security}

The paper proves EUF-CMA security, in the random oracle model, by
reduction to the \emph{discrete logarithm problem} in $\G$ ---
deliberately not to one-more discrete log, so that the Drijvers et al.\
metareduction does not apply. The proof is for an intermediate scheme,
\emph{FROST-Interactive}, in which the binding value is produced by a
one-time verifiable random function
$\rho_i = a_{ij} + b_{ij} \cdot H_\rho(m, B)$ with committed keys.

\begin{quotation}
\noindent\textbf{Theorem 6.1} (ePrint 2020/852). \emph{If the discrete
logarithm problem in $\G$ is $(\tau', \epsilon')$-hard, then the
FROST-Interactive signature scheme over $\G$ with $n$ signing
participants, a threshold of $t$, and a preprocess batch size of $\pi$ is
$(\tau, n_h, n_p, n_s, \epsilon)$-secure whenever}
\[
\epsilon' \le \frac{\epsilon^2}{2 n_h + (\pi + 1) n_p + 1}
\quad \text{and}
\]
\[
\tau' = 4\tau + \bigl(30 \pi n_p + (4t - 2) n_s + (n + t - 1) t + 6\bigr)
\cdot t_{\mathrm{exp}} + O(\pi n_p + n_s + n_h + 1),
\]
\emph{such that $t_{\mathrm{exp}}$ is the time of an exponentiation in
$\G$, assuming the number of participants compromised by the adversary is
less than the threshold $t$.}
\end{quotation}

Here $n_h$, $n_p$, $n_s$ bound the forger's random-oracle, preprocess,
and signing queries. The architecture is Bellare--Neven generalised
forking: the DL challenge $\omega$ is embedded as the honest
participant's free-coefficient commitment $\phi_{t0}$ in a simulated
KeyGen, the simulator is run twice on the same tape with oracle answers
diverging from index $J$, and two forgeries yield
$\mathrm{dlog}(Y) = (z' - z)/(h'_J - h_J)$. The VRF is what lets the
simulator win the programming race: it can compute $R$ before the forger
does and program $H_2(R, Y, m)$ with success probability $1$.

The step from FROST-Interactive to plain FROST (replace the VRF by
$\rho_i = H_1(i, m, B)$) is argued \emph{heuristically} in the paper's
\S6.2, not proven --- the Equation-1 gap of \S\ref{sec:fr-binding}. This
gap is what the later literature closed, differently and under a
different assumption (\S\ref{sec:fr-later}).

\subsection{The corrected and extended analyses}\label{sec:fr-later}

Three later results matter for how this volume cites FROST's security.

\paragraph{Crites--Komlo--Maller (ePrint 2021/1375).} This work, ``How
to Prove Schnorr Assuming Schnorr: Security of Multi- and Threshold
Signatures'' (IACR ePrint 2021/1375, received 2021-10-12, last revised
2022-08-03), introduced the optimisation now called \emph{FROST2}: a
single binding factor $d = h_1(vk, lr)$ shared by all signers,
replacing the original per-signer $d_i = h_1(vk, lr, i)$ (that original
is retroactively \emph{FROST1}) and reducing the signers' scalar
multiplications from linear in the number of signers to constant.

\paragraph{Bellare--Crites--Komlo--Maller--Tessaro--Zhu (CRYPTO 2022).}
The paper ``Better than Advertised Security for Non-interactive
Threshold Signatures'' --- its FROST/BLS part maintained as
Bellare--Tessaro--Zhu, ePrint 2022/833, June 2022 --- gave the analysis
now cited by RFC 9591. The authors define a hierarchy
$\mathrm{TS\text{-}UF\text{-}0} \le \cdots \le
\mathrm{TS\text{-}UF\text{-}4}$ (and strong variants
$\mathrm{TS\text{-}SUF\text{-}2/3/4}$) grading what counts as a trivial
forgery. Their results, proven in the ROM under the \emph{one-more
discrete logarithm} (OMDL) assumption, without the AGM, and in the
trusted-dealer setting (the DKG is explicitly not modelled):
\begin{itemize}
\item FROST1 is TS-SUF-3-secure, and is \emph{not} TS-UF-4-secure;
\item FROST2 is TS-SUF-2-secure, and is \emph{not} TS-UF-3-secure ---
  the shared binding factor loses the guarantee that the signer set that
  committed is the signer set that signed;
\item concretely (their Theorem 5.1, verbatim bound), for a TS-SUF-2
  adversary $A$ making at most $q_s$ preprocessing and $q_h$
  random-oracle queries there is an OMDL adversary $B$ making at most
  $2 q_s + \mathit{ns}$ challenge queries ($\mathit{ns}$ their number of
  parties) with
  \[
  \mathrm{Adv}^{\mathrm{ts\text{-}suf\text{-}2}}_{\mathrm{FROST2}[\G]}(A)
  \le \sqrt{q \cdot \bigl(\mathrm{Adv}^{\mathrm{omdl}}_{\G}(B)
  + 3 q^2 / p \bigr)}, \qquad q = q_s + q_h + 1,
  \]
  and an analogous Theorem 5.4 for FROST1 at TS-SUF-3.
\end{itemize}
The security FROST actually deploys with is therefore: OMDL + ROM, static
corruptions below $t$, trusted-dealer keys --- a different basis than the
original paper's DL-based proof of the interactive variant plus heuristic
bridge.

\paragraph{Gouv\^ea--Komlo (ePrint 2024/436).} The preprint
``Re-Randomized FROST'' (received 2024-03-13) is the design ZIP 312
builds on: FROST unchanged, plus algorithms $\mathsf{GenRand}$
($\alpha \xleftarrow{\$} \Z_p$),
$\mathsf{RandShare}$ ($\bar{x}_i \gets x_i + \hat{\alpha}$ with
$\hat{\alpha} = H_r(\alpha, \mathrm{aux})$, $\mathrm{aux}$ an optional
external-protocol transcript contribution), $\mathsf{RandPKShare}$
($\overline{PK}_i = PK_i \cdot g^{\hat{\alpha}}$) and $\mathsf{RandPK}$
($\overline{PK} = PK \cdot g^{\hat{\alpha}}$). Correctness rests on the
$\sum_{i \in S} \lambda_i = 1$ identity, giving
$z = r + c(x + \hat{\alpha})$. Its Theorem 1 proves unforgeability in the
ROM under \emph{algebraic} OMDL (AOMDL) --- the same assumptions as plain
FROST per BCKMTZ --- with the bound (its Equation 3)
\[
\mathrm{Adv}^{\mathrm{sec}}_{TS,A} \le
\sqrt{\, q \cdot \mathrm{Adv}^{\mathrm{aomdl}}_{D}(\lambda)
+ 2 q^2 / p - \mathrm{negl}(\lambda) \,},
\qquad q = q_h + q_s,
\]
%% Radical span verified against the PDF page render (p9-09.png, paper
%% p. 9, Eq. 3): it covers all three terms, including "- negl(lambda)".
via the local forking lemma; the randomizer being public lets the
reduction strip it ($z_0 \gets z^* - c^* H_r(\alpha^*, \mathrm{aux}^*)$)
during extraction. The coordinator chooses the randomizer and is trusted
for privacy and unlinkability only, not for unforgeability.
The full treatment of this paper belongs to the re-randomisation chapter
(\S\ref{sec:fr-rerand}).

\subsection{RFC 9591 against the paper: what the standard changed}
\label{sec:fr-rfc}

RFC 9591 (IRTF CFRG, Informational, June 2024; Connolly, Komlo, Goldberg,
Wood) is FROST as implementations ship it, including the
\texttt{frost-core} 3.0.0 workspace this volume cites. Reconciling it
against ePrint 2020/852:

\begin{itemize}
\item \emph{Two rounds only.} The RFC specifies the two-round protocol;
  the paper's single-round-with-preprocessing variant (batch parameter
  $\pi$, commitment server) is declared out of scope.
\item \emph{No DKG.} The paper's Figure-1 KeyGen is dropped; key
  generation is out of scope, with \emph{trusted-dealer} generation given
  in Appendix C only. The BCKMTZ analysis the RFC cites matches this
  framing (it, too, models a trusted dealer).
\item \emph{Binding factor inputs widened.} The paper binds
  $\rho_i = H_1(i, m, B)$. The RFC computes, per participant,
  \[
  \mathtt{binding\_factor}_i = H_1(\mathtt{PK\_enc} \,\|\,
  H_4(\mathtt{msg}) \,\|\, H_5(\mathtt{encoded\ commitment\ list}) \,\|\,
  \mathtt{enc}(i)),
  \]
  adding the \emph{group public key} to the hash and pre-hashing the
  message and commitment list. The per-participant structure means RFC
  9591 standardises FROST1: the RFC's \S7.2 rules the
  single-binding-factor (FROST2) optimisation NOT RECOMMENDED, citing the
  loss of the started-set-equals-signed-set guarantee (the TS-SUF-3 vs
  TS-SUF-2 separation of \S\ref{sec:fr-later}).
\item \emph{Hedged nonces.} The paper samples $(d, e)$ uniformly; the RFC
  hedges: each nonce is $H_3(r \,\|\, \mathtt{enc}(sk_i))$ with $r$ a
  fresh $32$-byte random string, so a weak RNG is backstopped by the key
  share. It warns that fully deterministic derivation enables complete
  key recovery in the multi-party setting.
\item \emph{Five domain-separated oracles.} $H_1$ (binding factor),
  $H_2$ (challenge), $H_3$ (nonce), $H_4$ (message pre-hash), $H_5$
  (commitment-list pre-hash), each with a per-ciphersuite
  \texttt{contextString}; the challenge input order
  $\mathtt{group\_comm\_enc}$ $\|$ $\mathtt{PK\_enc}$ $\|$
  $\mathtt{msg}$ equals the paper's $H_2(R, Y, m)$.
\item \emph{Roles renamed and channels stated.} The signature aggregator
  becomes the \emph{Coordinator}; the channel is assumed authenticated.
  The Coordinator SHOULD verify the aggregate signature and MAY verify
  individual shares --- the paper's per-share check
  $g^{z_i} = R_i \cdot Y_i^{c \lambda_i}$ survives as
  $\mathsf{verify\_signature\_share}$ --- with identifiable abort, as in
  the paper, and no robustness.
\item \emph{Ciphersuites.} The RFC ships FROST(Ed25519, SHA-512),
  FROST(ristretto255, SHA-512) (RECOMMENDED), FROST(Ed448, SHAKE256),
  FROST(P-256, SHA-256), and FROST(secp256k1, SHA-256); no Zcash suite
  appears in the standard itself. The two Zcash suites,
  FROST(Jubjub, BLAKE2b-512) and FROST(Pallas, BLAKE2b-512), are
  instead defined by ZIP 312, whose
  $H_2$ personalisation $\mathtt{Zcash\_RedPallasH}$ makes the FROST
  challenge equal RedDSA's $\mathsf{H}^{\circledast}$ --- the precise
  sense in which the standardised protocol plugs into the
  spend-authorisation verifier of the \emph{Ironwood Guide},
  \S``Authorisation of the spend: RedPallas''.
\item \emph{Security claim.} The RFC claims EUF-CMA per BCKMTZ, citing
  both the original paper and BCKMTZ --- i.e., the deployed claim rests
  on the OMDL/ROM analysis, not the paper's DL-based theorem for the
  interactive variant.
\end{itemize}

Classification: FROST as in RFC 9591 and its Zcash ciphersuites as in
ZIP 312 are \emph{specified}; ZIP 312 itself remains Draft status (zips
commit 6961098410, 2026-07-09), with its key-generation procedure
explicitly out of scope and PedPoP DKG for Zcash deployment
\emph{designed-but-unspecified} (implemented in \texttt{frost-core}'s
\texttt{dkg} module but not standardised in any ZIP or RFC).
%% Bin confirmed 2026-07-15: no ZIP or RFC specifies any DKG for Zcash
%% custody; zips PR #895, which would specify key generation, remains
%% open (21 commits, awaiting review by daira and nuttycom).

%% FROST Guide draft: Rerandomized FROST for RedPallas --- ePrint 2024/436,
%% ZIP 312, the deployed crates, and the custody surfaces.
%% Body-only LaTeX; labels sec:fr-*; register per ironwood-guide; 80 cols.
%%
%% Sources (all verified firsthand 2026-07-15):
%%   - C. P. L. Gouvea, C. Komlo (both Zcash Foundation), "Re-Randomized
%%     FROST", IACR ePrint 2024/436, https://eprint.iacr.org/2024/436
%%     (received 2024-03-13, approved 2024-03-15; Preprint, no venue;
%%     CC BY 4.0; Komlo footnote: "This work was not part of my University
%%     of Waterloo duties."). Read in full. Direct PDF fetch from
%%     eprint.iacr.org is blocked by a Cloudflare challenge (two hosts +
%%     in-browser); PDF obtained from Wayback snapshot 20260218073940 of
%%     https://eprint.iacr.org/2024/436.pdf, byte-identical to
%%     .wip/frost-drafts/sources/rerand-2024-436.pdf. Local copies
%%     eprint-2024-436.pdf/.txt; Fig. 5 and Eqs. (2)/(3) verified against
%%     the page renders p8-08.png, p9-09.png, p10-10.png.
%%   - RFC 9591 (base FROST): .wip/frost-drafts/rfc9591.txt.
%%   - zcash/zips at commit 6961098410 (2026-07-09): zip-0312.rst (the only
%%     FROST ZIP --- ZIP 313 is the obsolete fee ZIP); zip-2005.md ("Usage
%%     with FROST"); zip-0316.rst lines 618-624; zip-0271.md lines 114-135;
%%     protocol/protocol.tex 4.15 (spendauthsig), 5.4.7 (concretereddsa),
%%     SURK-CMA at lines 4989-5011.
%%   - ZcashFoundation/frost at 2016e44ba4 (2026-04-23, workspace 3.0.0):
%%     frost-rerandomized/src/lib.rs, frost-core/src/round1.rs.
%%   - ZcashFoundation/reddsa at 3792daa95e (2026-05-22, v0.5.2):
%%     src/frost/redpallas.rs, src/frost/redpallas/rerandomized.rs +
%%     README.md, src/orchard.rs, src/hash.rs.

\section{Re-Randomized FROST for RedPallas}\label{sec:fr-rerand}

The FROST of \S\ref{sec:fr-frost}, standardised as RFC 9591
(\S\ref{sec:fr-rfc}), signs for a \emph{fixed} group public key: the
aggregate $(R, z)$ verifies under the single-signer Schnorr equation for
$PK$, and under nothing else. Zcash spend authorisation verifies under no
fixed key at all. Every Sapling Spend and every Orchard Action carries a
freshly randomised validating key, and consensus checks the
spend-authorisation signature against that per-Action key --- a key which,
as \S\ref{sec:fr-rerand-gap} makes precise, no set of FROST shares
interpolates to unless every share is shifted. This chapter presents the
design that closes the gap: the Re-Randomized FROST of Gouv\^ea and Komlo
(IACR ePrint 2024/436), its Zcash specification in ZIP 312, the
ciphersuite FROST(Pallas, BLAKE2b-512), the deployed implementation in the
\texttt{frost-rerandomized} and \texttt{reddsa} crates, and the custody
surfaces the construction serves.

A rigor note before the material. RFC 9591 is the only artefact here at
standards rigor. The Gouv\^ea--Komlo paper is a preprint --- received
2024-03-13, no peer-reviewed venue --- and ZIP 312 carries Draft status
(zips commit 6961098410, 2026-07-09). Following the series rule, every
claim beyond the RFC is binned as \emph{specified}, \emph{designed but
unspecified}, or an \emph{open problem}; the closing classification is in
\S\ref{sec:fr-rerand-custody}.

\subsection{The mismatch: consensus verifies only randomised keys}
\label{sec:fr-rerand-gap}

The protocol specification (\S4.15, spend authorization signatures)
prescribes, for every Sapling Spend and every Orchard Action, the
following signer-side flow. The signer chooses a fresh spend-auth
randomizer $\alpha$ and:
\begin{enumerate}
\item samples $\alpha \gets \mathsf{SpendAuthSig.GenRandom}()$;
\item derives $\rsk =
  \mathsf{SpendAuthSig.RandomizePrivate}(\alpha, \ask)$;
\item derives $\rk =
  \mathsf{SpendAuthSig.DerivePublic}(\rsk)$;
\item proves the Spend or Action statement with $\alpha$ in the auxiliary
  (secret) input and $\rk$ in the primary input;
\item signs: $\mathsf{spendAuthSig} =
  \mathsf{SpendAuthSig.Sign}_{\rsk}(\mathsf{SigHash})$.
\end{enumerate}
Consensus therefore verifies the signature under
$\rk = \aksym + [\alpha]\,G^{\mathsf{SpendAuth}}$, never
under $\aksym$ itself.

The randomisation algorithms are RedDSA's, from the protocol
specification (\S5.4.7):
\begin{align*}
\mathsf{RedDSA.RandomizePrivate}(\alpha, \mathsf{sk})
  &:= \mathsf{sk} + \alpha \pmod{r_{\G}}, \\
\mathsf{RedDSA.RandomizePublic}(\alpha, \mathsf{vk})
  &:= \mathsf{vk} + [\alpha]\,G,
\end{align*}
with identity randomizer $0$; $\mathsf{RedDSA.GenRandom}$ chooses $T$
uniformly among byte sequences of length $(\ell_H + 128)/8$ --- $80$
bytes at $\ell_H = 512$ --- and returns $H^{\circledast}(T)$, where
$H^{\circledast}(B) := \mathrm{LEOS2IP}_{\ell_H}(H(B)) \bmod r_{\G}$.
Signing draws its nonce as
$r = H^{\circledast}(T \,\|\, \mathsf{vk}^\star \,\|\, M)$ for fresh
uniform $T$, sets $R = [r]\,G$ and
$S = \bigl(r + H^{\circledast}(R^\star \,\|\, \mathsf{vk}^\star \,\|\, M)
\cdot \mathsf{sk}\bigr) \bmod r_{\G}$, and outputs
$\sigma = R^\star \,\|\, S^\star$. Validation recomputes
$c = H^{\circledast}(R^\star \,\|\, \mathsf{vk}^\star \,\|\, M)$ and
accepts iff $R \ne \bot$, $S < r_{\G}$, $R^\star$ is canonical
(post-ZIP-216), and
$[h_{\G}]\bigl(-[S]\,G + R + [c]\,\mathsf{vk}\bigr) = 0_{\G}$. RedPallas
specialises the template with $\G$ the Pallas group, $\ell_H = 512$, and
$H(x) = \mathrm{BLAKE2b}\text{-}512(\texttt{Zcash\_RedPallasH}, x)$;
$\mathsf{SpendAuthSig}^{\mathsf{Orchard}}$ uses the generator
$G^{\mathsf{SpendAuth}} =
\mathsf{GroupHash}^{\Pallas}(\mathtt{z.cash:Orchard}, \texttt{"G"})$ and
``is instantiated as RedPallas with key re-randomization''. The security
notion the protocol requires of the scheme is SURK-CMA (Strong
Unforgeability with Re-randomized Keys), adapted from Fleischhacker et
al.\ (PKC 2016, \S3), whose oracle answers queries of the form
$(m, \alpha)$ --- message and adversary-chosen randomizer together.

The construction and its single-party security are settled in the
\emph{Ironwood Guide}: \S``Authorisation of the spend: RedPallas'' builds
$\rsk := \ask + \alpha$ and
$\rk := \aksym + [\alpha]\,G^{\mathsf{SpendAuth}}$, and
\S``RedPallas unforgeability and re-randomisation'' proves EUF-CMA under
DL in the ROM \emph{including} adversary-chosen $\alpha$, by key-prefixed
forking with $\ask = \rsk^* - \alpha^*$. Per the layering
rule we cite those results and do not re-prove them.

The threshold problem is now visible. A FROST sharing of $\ask$
gives shares $\mathsf{sk}_i$ with
$\sum_{i \in S} \lambda_i \mathsf{sk}_i = \ask$ for every signing
set $S$; every signature that coalition can produce verifies under
$\aksym$. But no honest transaction ever asks for a signature under
$\aksym$: the verifier's key is $\rk$, shifted by a fresh
$\alpha$ per Action. The shares must be shifted too --- and, as the next
subsection shows, the correct shift is the \emph{same} $\alpha$ for every
participant, not a Shamir sharing of $\alpha$.

\subsection{Re-Randomized FROST: the construction}
\label{sec:fr-rerand-scheme}

The design is due to Gouv\^ea and Komlo (both Zcash Foundation),
``Re-Randomized FROST'', IACR ePrint 2024/436, received 2024-03-13; we
work from the full text, with the construction figure and both displayed
equations verified against page renders of the PDF. As in
\S\ref{sec:fr-frost} we quote the paper in its own multiplicative
notation ($g^x$ for $[x]\,G$, $PK \cdot g^{\hat{\alpha}}$ for
$PK + [\hat{\alpha}]\,G$); the additive RedPallas form returns with
ZIP 312 in \S\ref{sec:fr-rerand-zip}.

\begin{definition}[Perfectly rerandomizable threshold scheme;
ePrint 2024/436, Definition 5]\label{def:fr-rerandomizable}
A threshold signature scheme
\[
TS = (\mathsf{Setup}, \mathsf{KeyGen}, (\mathsf{Sign}, \mathsf{Sign}'),
\mathsf{Combine}, \mathsf{Verify})
\]
is \emph{perfectly re-randomizable}
if there exist additional PPT algorithms $\mathsf{GenRand}$,
$\mathsf{RandShare}$, $\mathsf{RandPKShare}$, and $\mathsf{RandPK}$,
and a randomness distribution $X$, such that the
$\mathsf{Rand}^*$ algorithms take a randomizer $\alpha \in X$ and an
auxiliary string $\mathrm{aux} \in \{0,1\}^*$.
\end{definition}

The paper's Remark 1 explains $\mathrm{aux}$: it models an optional
contribution from an external protocol transcript, so that a new session
changes the re-randomised key even under a repeated $\alpha$.

\begin{construction}[Re-Randomized FROST; ePrint 2024/436, Figure 5]
\label{def:fr-rerand-frost}
The additions to plain FROST --- the paper's grey-boxed lines --- are
exactly four algorithms:
\begin{align*}
\mathsf{GenRand}() &: \alpha \xleftarrow{\$} \Z_p; \\
\mathsf{RandShare}(x_i, \alpha, \mathrm{aux}) &:
  \hat{\alpha} \gets H_r(\alpha, \mathrm{aux}); \quad
  \bar{x}_i \gets x_i + \hat{\alpha}; \\
\mathsf{RandPKShare}(PK_i, \alpha, \mathrm{aux}) &:
  \overline{PK}_i \gets PK_i \cdot g^{\hat{\alpha}}; \\
\mathsf{RandPK}(PK, \alpha, \mathrm{aux}) &:
  \overline{PK} \gets PK \cdot g^{\hat{\alpha}}.
\end{align*}
The FROST core runs unchanged on the shifted inputs.
Algorithm $\mathsf{Sign}()$ samples $(r_k, s_k) \xleftarrow{\$} \Z_p^2$
and commits $(R_k, S_k) = (g^{r_k}, g^{s_k})$. Algorithm
$\mathsf{Sign}'$, given the
randomised share $\bar{x}_k$ and randomised key $\overline{PK}$ in place
of $x_k$ and $PK$, derives each binding factor $a_i$ by $H_{\mathrm{non}}$
from the participant identifier, the message, $\overline{PK}$, and the
full commitment list $\{(i, R_i, S_i)\}_{i \in S}$, then computes
\[
R \gets \prod_{i \in S} R_i \cdot S_i^{a_i}, \qquad
c \gets H_{\mathrm{sig}}(\overline{PK}, R, m), \qquad
z_k \gets r_k + s_k a_k + c\,\lambda_k \bar{x}_k .
\]
Algorithm $\mathsf{Combine}$ outputs $z \gets \sum_{i \in S} z_i$,
$\sigma \gets (R, z)$; $\mathsf{Verify}(\overline{PK}, m, \sigma)$
recomputes the challenge and accepts iff
$R \cdot \overline{PK}^{\,c} = g^z$.\footnote{The paper's own argument
order for $H_{\mathrm{sig}}$ is inconsistent: $\mathsf{Sign}'$ and
$\mathsf{Combine}$ write $H_{\mathrm{sig}}(PK, R, m)$, while
$\mathsf{Verify}$ and the reduction's table write
$H_{\mathrm{sig}}(PK, m, R)$ --- and both differ from the order RFC 9591
(\S4.6) and RedDSA fix,
$c = H_2(R^\star \,\|\, PK^\star \,\|\, m)$, which ZIP 312 and the
deployed \texttt{reddsa} crate use. We read the paper's orderings as a
presentational slip with no normative weight; the deployed order is
RedDSA's. The paper also twice refers to ``Figure 4'' where Figure 5 is
meant.}
\end{construction}

The paper is explicit that this is the whole design: ``We do not make any
changes to the plain FROST protocol.'' The only interface change is which
share and which public key $\mathsf{Sign}'$ is handed.

\emph{Correctness} is one identity. The paper's Equation (2):
\begin{equation}\label{eq:fr-rerand-z}
z = \sum_{i \in S} r_i + \sum_{i \in S} s_i a_i
    + c \sum_{i \in S} \lambda_i \bar{x}_i
  = \sum_{i \in S} r_i + \sum_{i \in S} s_i a_i
    + c \Bigl( \sum_{i \in S} \lambda_i x_i
    + \hat{\alpha} \sum_{i \in S} \lambda_i \Bigr)
  = r + c\,(x + \hat{\alpha}) = r + c\,\hat{x},
\end{equation}
``because $x$ is Shamir secret shared'' and
``$\sum_{i \in S} \lambda_i = 1$ since it is the interpolation of the
constant polynomial $f(X) = 1$''. The output is a regular re-randomised
signature for $\overline{PK} = PK \cdot g^{\hat{\alpha}}$ with
$\hat{\alpha} = H_r(\alpha, \mathrm{aux})$.

The identity deserves emphasis, because it is the entire reason the
construction is this simple. Every participant adds the \emph{same}
$\hat{\alpha}$ to its share --- no dealing, no per-participant shares of
the randomizer --- and the Lagrange weights, summing to one over any
interpolating set, carry the common shift through interpolation so that
the group secret moves by $\hat{\alpha}$ rather than by
$|S| \cdot \hat{\alpha}$. Had the weights summed to anything else, or
varied with $S$ in a way the signers could not predict, the shifted
shares would interpolate to a key nobody computed.

\subsection{Unforgeability: TRUF, AOMDL, and Theorem 1}
\label{sec:fr-rerand-sec}

\paragraph{The game.} The paper's unforgeability notion (Definition 6,
game $\mathrm{TRUF}$, its Figure 4) extends threshold unforgeability with
adversarial randomizers. Corruption is \emph{static}: the adversary fixes
its up-to-$(t-1)$ corrupt parties before key generation (the paper titles
its Figures 2 and 4 ``static unforgeability games''). The signing oracle
$\mathcal{O}^{\mathsf{Sign}'}(k, \mathrm{ssid}, m, S, \alpha,
\mathrm{aux}, \{\rho_i\})$ lets the \emph{adversary} choose the
randomizer and the auxiliary string per query; a forgery
$(m^*, \sigma^*, \alpha^*, \mathrm{aux}^*)$ wins if $m^*$ was never
queried and $\sigma^*$ verifies under $PK$ \emph{or} under
$\mathsf{RandPK}(PK, \alpha^*, \mathrm{aux}^*)$. The unrandomised notion
is the special case randomizer $= 0$. The lineage is Fleischhacker et
al.\ (PKC 2016) --- the same single-party re-randomisable-key framework,
and re-randomised-Schnorr EUF result, that the protocol specification's
SURK-CMA definition adapts (\S\ref{sec:fr-rerand-gap}).

\paragraph{The assumption.} AOMDL is the \emph{algebraic one-more
discrete logarithm} game (the paper's Figure 1, taken from the MuSig2
paper): the adversary receives $\ell + 1$ challenges, may query a
discrete-log oracle only on linear combinations
$(\alpha, \{\beta_i\})$ of environment-generated elements, and must
return all $\ell + 1$ discrete logarithms having made at most $\ell$
queries. The algebraic restriction on oracle queries is what
\S\ref{sec:fr-later}'s OMDL does not impose.

\paragraph{The theorem.} Verbatim (Theorem 1 with Equation (3); the
radical verified against the rendered page):

\begin{quotation}
\noindent\textbf{Theorem 1} (ePrint 2024/436). \emph{Rerandomized-FROST
is unforgeable in the random oracle model, assuming the AOMDL assumption
holds in $\G$}, concretely
\begin{equation}\label{eq:fr-rerand-bound}
\mathrm{Adv}^{\mathrm{sec}}_{TS,\mathcal{A}} \le
\sqrt{\, q \cdot \mathrm{Adv}^{\mathrm{aomdl}}_{D}(\lambda)
      + 2 q^2 / p - \mathrm{negl}(\lambda) \,},
\qquad q = q_h + q_s,
\end{equation}
with $q_h$ random-oracle and $q_s$ signing queries.
\end{quotation}

\begin{remark}[Placement of the negligible term]\label{rem:fr-negl}
The printed bound places $-\,\mathrm{negl}(\lambda)$ \emph{inside} the
square root; we transcribe it as printed. The customary form of such
forking-lemma bounds --- including the BCKMTZ bound quoted in
\S\ref{sec:fr-later} --- carries additive slack outside the radical, and
a negative term inside is dimensionally odd for an advantage bound.
\end{remark}

\paragraph{Proof shape.} The reduction $\mathcal{B}$ receives
$2 q_s + 1$ AOMDL challenges and sets $PK = X_0$. It simulates
trusted-dealer key generation by sampling the corrupt parties' shares
and producing every verification share by Lagrange interpolation in the
exponent, so that $PK_i$ is implicitly $g^{f(i)}$ without $\mathcal{B}$
knowing any honest share. Round-one answers are challenge pairs
$R_i = X_i$, $S_i = X_{i+1}$. In round two the reduction runs
$\mathsf{RandPK}$ and $\mathsf{RandPKShare}$ honestly --- it cannot run
$\mathsf{RandShare}$, holding no share --- and derives each honest
signature share in the exponent, as
$Z_k = R_i \cdot S_i^{a_k} \cdot \widehat{PK}_k^{\,c}$, querying the
AOMDL linear-combination oracle for the scalar $z_k$; the simulation is
perfect. The local forking lemma of Bellare--Dai--Li (ASIACRYPT 2019)
reprograms $H_{\mathrm{sig}}$ at the forgery's challenge query. The
randomizers, being public, are stripped before extraction: when
$\alpha \ne \bot$,
\[
z_0 = z^{*} - c^{*} H_r(\alpha^{*}, \mathrm{aux}^{*}), \qquad
z_1 = z^{**} - c^{**} H_r(\alpha^{**}, \mathrm{aux}^{**}), \qquad
y_0 = \frac{z_0 - z_1}{c^{*} - c^{**}},
\]
and the remaining challenge logarithms are extracted as in the FROST
proof of Bellare--Crites--Komlo--Maller--Tessaro--Zhu (CRYPTO 2022;
\S\ref{sec:fr-later}). The count closes the one-more game: $2 q_s$
oracle queries against $2 q_s + 1$ extracted logarithms.

\paragraph{What the theorem does not say.} Four caveats frame every
citation of this result. First, the paper proves \emph{unforgeability
only}: it contains no formal unlinkability definition and no
unlinkability theorem; that property is argued at the ZIP layer
(\S\ref{sec:fr-rerand-zip}) by matching RedDSA's distributions, under
the protocol specification's SURK-CMA framing. Second, corruptions are
static. Third, the threshold-scheme definition assumes centralised
(trusted-dealer) key generation, which the paper says ``can [be]
adapted'' to a DKG --- adapted, not proven. Fourth, the paper's \S5.1
sketches how a DKG-style contribution round could remove the
trusted-for-privacy coordinator (all parties contribute to the
randomizer, none learns it) at the cost of extra communication rounds;
ZIP 312 does not adopt this. The cost of re-randomisation itself is one
fixed-base scalar multiplication per participant and per coordinator;
the paper's Table 1 benchmarks (Pallas suite, Ryzen 9 5900X) put the
signing overhead at $83\%$ with $2$ signers, $45\%$ at $7$, and $8\%$
at $67$.

\subsection{ZIP 312: the specification}\label{sec:fr-rerand-zip}

ZIP 312, ``FROST for Spend Authorization Multisignatures'' (Status
Draft, Category Wallet, created 2022-08; owners Gouv\^ea, Komlo,
Connolly; zcash/zips PR \#662), specifies Re-Randomized FROST as a
delta against RFC 9591. Its target property beyond unforgeability is
\emph{unlinkability}, which the ZIP defines as statistical independence
of signatures from the signers' long-term keys: for uniform randomizers
and absent metadata leakage, it is ``impossible to determine whether two
transactions were generated by the same party''.

\paragraph{The modifications to RFC 9591.} Five, and only five:
\begin{itemize}
\item \emph{Randomizer generation.} A new helper
  \texttt{randomizer\_generate(msg, commitment\_list)}: draw
  \texttt{rng\_randomizer} $\gets$ \texttt{random\_bytes(Ns)}, form
  \begin{align*}
  \texttt{randomizer\_input} = {} & \texttt{rng\_randomizer} \,\|\, \\
  & \texttt{encode\_signing\_package(commitment\_list, msg)},
  \end{align*}
  and return $H_R(\texttt{randomizer\_input})$. Implementations MAY use
  another encoding ``as long as all values (the message, and the
  identifier, binding nonce and hiding nonce for each participant) are
  unambiguously encoded''.
\item \emph{Round one} is unchanged.
\item \emph{Round two.} The Coordinator generates the randomizer and
  sends it to each signer ``over a confidential and authenticated
  channel'' --- plain FROST requires authentication only --- together
  with the message and the commitment list. Each participant then runs
  \texttt{sign} with $\mathsf{sk}_i + \mathrm{randomizer}$ in place of
  its share and $PK + [\mathrm{randomizer}]\,G$ in place of the group
  public key.
\item \emph{Aggregation} shifts the group key the same way:
  $PK \mathrel{+}= [\mathrm{randomizer}]\,G$.
\item \emph{Share verification} shifts both keys:
  $PK_i \mathrel{+}= [\mathrm{randomizer}]\,G$ and
  $PK \mathrel{+}= [\mathrm{randomizer}]\,G$ in
  \texttt{verify\_signature\_share}.
\end{itemize}
The ZIP's ``randomizer'' is the paper's \emph{effective} randomizer
$\hat{\alpha}$: per the paper's \S7 implementation notes,
$\mathsf{GenRand}$ and $H_r$ are folded into the single hash call above,
and the value distributed to signers is the already-hashed scalar.

\paragraph{The algebra, in RedDSA notation.} The ZIP's Rationale
restates the correctness identity of \eqref{eq:fr-rerand-z} additively,
and it is worth displaying because it is the entire consensus-facing
content of the design. A valid RedDSA response must satisfy
$z = r + c \cdot \mathsf{sk}$. Each FROST share contributes
\[
z_i = \bigl(\mathtt{hiding\_nonce}_i
      + \mathtt{binding\_nonce}_i \cdot \rho_i\bigr)
      + \lambda_i \cdot c \cdot \mathsf{sk}_i ,
\]
and under re-randomisation the key terms become
$c \cdot (\mathsf{sk} + \mathrm{randomizer})$ on the single-signer side
and, summed over the coalition,
\[
\sum_{i \in S} \lambda_i \, c \, (\mathsf{sk}_i + \mathrm{randomizer})
= c \Bigl( \sum_{i \in S} \lambda_i \mathsf{sk}_i
  + \mathrm{randomizer} \cdot \sum_{i \in S} \lambda_i \Bigr)
= c \, (\mathsf{sk} + \mathrm{randomizer}),
\]
since $\sum_i \lambda_i \mathsf{sk}_i = \mathsf{sk}$ by Shamir and
$\sum_i \lambda_i = 1$ (the ZIP proves the latter in its
\texttt{sum-lambda-proof} footnote, citing math.SE 1325342). Every
participant adds the same randomizer to its share; the
weights-sum-to-one identity is exactly why the group secret shifts by
$\alpha$ and not by $|S| \cdot \alpha$.

\paragraph{Who learns the randomizer.} The threat model is the ZIP's
sharpest departure from RFC 9591. The Coordinator generates the
randomizer and is ``trusted with the privacy of the transaction (which
includes the unlinkability property)'' but \emph{not} with
unforgeability: ``a rogue Coordinator \ldots\ should not be able to
create signed transactions without the approval of
\texttt{MIN\_PARTICIPANTS} participants''. All key-share holders receive
the randomizer and are likewise trusted with privacy only. The role
assignment is natural: the Coordinator fits the transaction builder, who
already creates the zero-knowledge proof --- and the proof needs
$\alpha$ as auxiliary input (step 4 of the flow in
\S\ref{sec:fr-rerand-gap}). The confidentiality requirement is not
optional hygiene: anyone learning $\hat{\alpha}$ can un-randomise the
on-chain key, $\aksym = \rk - [\hat{\alpha}]\,G$, so the
paper advises encrypting the randomizer ``so that eavesdroppers can't
break unlinkability''. Non-requirements are stated with equal care: the
coordinator-less variant of RFC 9591 \S7.5 is unsupported;\footnote{ZIP
312's reference for the variant mislabels it as RFC 9591 ``Section 7.3:
Removing the Coordinator Role''; the section is 7.5 (\S7.3 is ``Nonce
Reuse Attacks''), and the reference's URL anchor targets the correct
section.} share holders
can always link the signing operation to the eventual on-chain
transaction; key generation is out of scope; network privacy is out of
scope.

\paragraph{The message is a hash.} The signed message is the
$\mathsf{SigHash}$, which conveys nothing about the transaction to a
signer reading it. Signers MUST therefore check the
$\mathsf{SigHash}$ against transaction data sent over the same channel,
or recompute it themselves; the mechanism is out of the ZIP's scope.
This is the threshold echo of a rule the single-signer flow gets for
free: a signer who cannot see what it signs authorises whatever the
builder chose.

\paragraph{Unlinkability, and where it is proven.} Nowhere, formally:
the paper proves unforgeability only (\S\ref{sec:fr-rerand-sec}), and
the ZIP argues unlinkability in its Rationale by matching distributions
--- ``\texttt{randomizer\_generate} generates \texttt{randomizer}
uniformly at random as required by \texttt{RedDSA.GenRandom}'', and the
signing flow is compatible with $\mathsf{RedDSA.RandomizePrivate}$,
$\mathsf{RedDSA.RandomizePublic}$, $\mathsf{RedDSA.Sign}$, and
$\mathsf{RedDSA.Validate}$ --- while its Requirements defer the
formal notion to the protocol specification's SURK-CMA definition. The
distribution-matching argument does connect to a proved statement: the
\emph{Ironwood Guide}'s \S``RedPallas unforgeability and
re-randomisation'' shows honest uniform re-randomisations are
distributed as fresh Schnorr keys. What no source supplies is a
threshold-specific unlinkability theorem. Bin: the mechanism is
\emph{specified} (ZIP 312, Draft); its unlinkability guarantee rests on
a distribution argument, not a theorem.

\subsection{The ciphersuite FROST(Pallas, BLAKE2b-512)}
\label{sec:fr-rerand-suite}

ZIP 312 defines the ciphersuite in RFC 9591's template. The group is
Pallas with base point $G^{\mathsf{SpendAuth}}$ as in protocol \S5.4.7.1,
of order $r_{\Pallas}$; \texttt{RandomScalar} is uniform in
$[0, \mathtt{Order} - 1]$; \texttt{SerializeElement} is the point
compression $\mathrm{repr}_{\Pallas}$ ($P^\star$ in series notation) and
\texttt{DeserializeElement} is $\mathrm{abst}_{\Pallas}$, erroring on
$\bot$ \emph{and on the identity element}; \texttt{SerializeScalar} is
little-endian over $32$ bytes and \texttt{DeserializeScalar} rejects
values $\ge$ the order. The hash is BLAKE2b-512 with $16$-byte
personalisation strings and $N_h = 64$; scalar-valued hashes interpret
the $64$ output bytes as a little-endian integer reduced mod the order.
Table~\ref{tab:fr-pallas-suite} lists the oracles.

\begin{table}[htbp]
\centering
\small
\begin{tabular}{@{}lll@{}}
\toprule
Oracle & Role & Personalisation \\
\midrule
$H_1$ & binding factor        & \texttt{FROST\_RedPallasR} \\
$H_2$ & challenge             & \texttt{Zcash\_RedPallasH} \\
$H_3$ & nonce                 & \texttt{FROST\_RedPallasN} \\
$H_4$ & message pre-hash      & \texttt{FROST\_RedPallasM} \\
$H_5$ & commitment-list pre-hash & \texttt{FROST\_RedPallasC} \\
$H_R$ & randomizer            & \texttt{FROST\_RedPallasA} \\
\bottomrule
\end{tabular}
\caption{FROST(Pallas, BLAKE2b-512) hash personalisations (ZIP 312).
The challenge hash reuses RedPallas's own personalisation; the FROST
additions use the \texttt{FROST\_} prefix.}
\label{tab:fr-pallas-suite}
\end{table}

The single load-bearing row is $H_2$: ZIP 312 states it is
``equivalent to $H^{\circledast}(m)$'' of RedPallas, so the FROST
challenge \emph{is} the on-chain challenge, and signature verification
for the suite is RedPallas validation itself. The ZIP's compatibility
rationale completes the argument that a FROST aggregate is a plain
RedPallas signature: the $(R, S)$ output split matches
$\mathsf{RedDSA.Validate}$'s parsing; FROST and Zcash use the same
challenge input order ($R$, public key, message); and the fact that $r$
and $R$ are FROST-generated rather than produced by
$\mathsf{RedDSA.Sign}$'s nonce derivation is irrelevant to a verifier,
which sees only the distribution of $R$. A parallel suite
FROST(Jubjub, BLAKE2b-512) --- personalisations
\texttt{FROST\_RedJubjub\{R,N,M,C,A\}} and \texttt{Zcash\_RedJubjubH}
--- serves Sapling spend authorisation; this volume works the Pallas
suite and treats the Jubjub twin as its mechanical transliteration.

\subsection{Deployed implementation: \texttt{frost-rerandomized} and
\texttt{reddsa}}\label{sec:fr-rerand-impl}

Claims in this subsection cite ZcashFoundation/frost at commit
\texttt{2016e44} (2026-04-23, workspace version 3.0.0) and
ZcashFoundation/reddsa at commit \texttt{3792daa} (2026-05-22, crate
version 0.5.2), examined 2026-07-15. The pinned commits are the
published releases --- tag \texttt{frost-core/v3.0.0} (crates.io
2026-04-23) and \texttt{reddsa} 0.5.2 --- so main-at-commit and the
released crates coincide. ZIP 312 names \texttt{reddsa} as the
reference implementation of both ciphersuites.

\paragraph{The generic layer.} Crate \texttt{frost-rerandomized}
(\texttt{frost-rerandomized/src/lib.rs}) implements
Construction~\ref{def:fr-rerand-frost} over any \texttt{frost-core}
ciphersuite. Trait impl \texttt{Randomize<KeyPackage>} computes the
randomised signing share as $\mathtt{signing\_share} + \mathtt{randomizer}$
and the randomised verifying share by adding the
\texttt{randomizer\_element}; struct \texttt{RandomizedParams} carries
the \texttt{randomizer} (``also called alpha'' in its documentation),
$\mathtt{randomizer\_element} = [\mathtt{randomizer}]\,G$, and
$\mathtt{randomized\_verifying\_key} = PK + \mathtt{randomizer\_element}$.
The current flow differs instructively from a literal reading of the
ZIP: the Coordinator sends a randomizer \emph{seed}, not the
randomizer, and each participant recomputes the randomizer from the
seed and its own round-one commitments. The derivation, its divergence
from the draft text --- the message is no longer a hash input --- and
the reconciliation in flight in zips PR \#895 are treated once, in
\S\ref{sec:fr-deploy-rand}.

\paragraph{The RedPallas suite.} Module \texttt{src/frost/redpallas.rs}
in \texttt{reddsa} implements the \texttt{frost-core}
\texttt{Ciphersuite} trait for \texttt{PallasBlake2b512} (contextual ID
\texttt{"FROST(Pallas, BLAKE2b-512)"}). Its \texttt{generator()} is
\texttt{orchard::SpendAuth::basepoint()}, whose constant bytes are
reproducible as
\texttt{pallas::Point::hash\_to\_curve("z.cash:Orchard")(b"G")} --- the
protocol's $G^{\mathsf{SpendAuth}}$. Oracles $H_1/H_3/H_4/H_5$ carry the
\texttt{FROST\_RedPallas\{R,N,M,C\}} personalisations of
Table~\ref{tab:fr-pallas-suite}; $H_2$ goes through the crate's
\texttt{HStar} default, \texttt{Zcash\_RedPallasH}; the
\texttt{RandomizedCiphersuite::hash\_randomizer} impl uses
\texttt{FROST\_RedPallasA} (the ZIP's $H_R$); scalar reduction is the
$64$-byte-wide \texttt{from\_uniform\_bytes} (\texttt{src/hash.rs}).
Beyond ZIP 312, the crate defines two further personalisations ---
\texttt{HDKG} $=$ \texttt{FROST\_RedPallasD} and \texttt{HID} $=$
\texttt{FROST\_RedPallasI}, the DKG and identifier hashes --- for key
generation, which the ZIP does not specify at all.

\paragraph{Even-$Y$ normalisation.} One deployment detail appears in
neither the paper nor the ZIP 312 body: the \texttt{reddsa} hooks
\texttt{post\_generate}/\texttt{post\_dkg} normalise fresh key material
to the even-$Y$ representative the protocol specification requires of
$\aksym$; the full mechanism is \S\ref{sec:fr-deploy-stack}. The
\emph{randomised} key $\rk$ need not satisfy the even-$Y$ condition ---
only $\aksym$ does.

\paragraph{The RedPallas re-randomised API.} Module
\texttt{src/frost/redpallas/rerandomized.rs} re-exports the
re-randomised flow for the suite ---
\texttt{sign\_with\_randomizer\_seed()}, \texttt{aggregate()}, and
\texttt{aggregate\_custom()} --- together with the type aliases
\texttt{Randomizer} and \texttt{RandomizedParams} over
\texttt{PallasBlake2b512}. The module README's worked example runs the
full trusted-dealer plus re-randomised signing flow and ends by
verifying the aggregate under the
\texttt{randomized\_verifying\_key()} of its \texttt{RandomizedParams}
--- the executable form of this chapter's claim that a FROST aggregate
is a plain RedPallas signature under $\rk$.

\subsection{Custody surfaces and classification}
\label{sec:fr-rerand-custody}

\paragraph{(a) Spend-authorisation custody.} The primary surface: a
$t$-of-$n$ sharing of $\ask$ signing Orchard Actions (and, via
the Jubjub twin suite, Sapling Spends). FROST replaces step 5 of the
signer flow in \S\ref{sec:fr-rerand-gap}; steps 1--4 move to the
Coordinator-as-builder, who draws the randomizer, proves the Action
statement with $\alpha$ as auxiliary input, and publishes $\rk$
--- the circuit and consensus rules are untouched, per the
\emph{Ironwood Guide}, \S``Authorisation of the spend: RedPallas''. Bin:
\emph{specified} (ZIP 312, Draft), with key generation explicitly out of
the ZIP's scope; DKG-based key generation for Zcash remains
\emph{designed but unspecified} (implemented in \texttt{frost-core},
standardised nowhere), as \S\ref{sec:fr-rfc} already concluded for the
base protocol. Related anticipations elsewhere in the ZIP corpus: ZIP
316 requires FROST unified viewing key material to follow ZIPs 312 and
2005, and ZIP 271 anticipates lockbox disbursement to ``a FROST
address'', noting (as of May 2025) that ZIP 312 does not specify key
generation.

\paragraph{(b) ZIP 2005 quantum spending key custody.} ZIP 2005's
``Usage with FROST'' section constrains threshold deployments of its
quantum-recovery hierarchy: the key \aksym{} is derived jointly, by
FROST DKG, from participants' shares of \ask, and the participants
privately agree on a value $\mathsf{sk}$ from which $\nksym$,
$\mathsf{qsk}$, $\mathsf{qk}$, and $\mathsf{rivk}_{\mathrm{ext}}$
derive. The full ceremony, its derivations, and its asymmetric threat
model --- every participant holds $\mathsf{qsk}$, so the threshold
property does not survive a quantum adversary --- are developed once,
in \S\ref{sec:fr-deploy-qsk}. Bin: the FROST-facing requirements are
\emph{specified} (ZIP 2005's requirements list calls for full
compatibility with ZIP 312 FROST, hardware wallets, and their
combination); the combined DKG-plus-shared-$\mathsf{sk}$ key generation
is \emph{designed but unspecified} --- no source specifies the protocol
by which participants ``privately agree'' on $\mathsf{sk}$.
%% Shared-sk bin confirmed 2026-07-15: zips PR #895 drafts a
%% contributory sk-generation section (commit 2026-03-05), still
%% unmerged; designed-but-unspecified stands, #895's merge the trigger.

\paragraph{(c) Crosslink finalizer keys.} The \emph{Crosslink Guide},
\S``Finalizer keys'', bins threshold custody of finalizer keys as an
\emph{open problem}, and we adopt that bin. One precision for this
volume: the applicable protocol there is \emph{plain} FROST --- the
finalizer design has no key re-randomisation, so the natural suite is
RFC 9591's FROST(Ed25519, SHA-512), not either ZIP 312 suite. The
prototype state and this observation are developed once, in
\S\ref{sec:fr-deploy-crosslink}.
%% TODO: confirm with the parent that the intended framing is
%% "candidate custody surface for plain FROST (RFC 9591 ed25519 suite)",
%% not rerandomized FROST.

\paragraph{(d) PCZT signing.} The multiparty transaction flow that
would host the Coordinator role is the PCZT of the \emph{Wallet Guide},
\S``Partially created transactions (PCZT)'': ZIP 374 (Partially Created
Zcash Transaction Format, Draft, owner Jack Grigg) names ``FROST
threshold signing (RFC 9591)'' in its Motivation among the multiparty
use cases its Signer role must accommodate.

\paragraph{Classification.} In the series' three bins: base FROST and
its five ciphersuites are \emph{specified} at standards rigor (RFC 9591,
\S\ref{sec:fr-rfc}); Re-Randomized FROST, its threat model, and both
Zcash ciphersuites are \emph{specified} by ZIP 312 at Draft status, with
the unforgeability theorem resting on a preprint (ePrint 2024/436, no
venue) and unlinkability on a distribution argument rather than a
theorem; key generation --- FROST DKG for Zcash, and ZIP 2005's
shared-$\mathsf{sk}$ agreement --- is \emph{designed but unspecified};
threshold custody of Crosslink finalizer keys is an \emph{open problem}.
Deployed-behaviour claims (seed-based randomizer derivation, even-$Y$
normalisation, the extra DKG personalisations) are implementation facts
of the pinned commits, cited as such.

%% FROST Guide draft: deployment surfaces and open questions.
%% Body-only LaTeX; labels sec:fr-deploy*. Register per ironwood-guide.
%%
%% Sources (all verified firsthand 2026-07-15):
%%   - ZcashFoundation/frost @ 2016e44ba4a4757a996300350063b937a2ad33e8
%%     (2026-04-23, tag frost-core/v3.0.0, workspace version 3.0.0).
%%   - ZcashFoundation/reddsa @ 3792daa95e588c1af6bd4805105bfb6ea7e9ad49
%%     (2026-05-22, tag 0.5.2).
%%   - zcash/zips @ 69610984109078075e7e64989ecf89d63a259b97 (2026-07-09):
%%     zip-0312.rst (Draft), zip-2005.md (Proposed), zip-0374.md.
%%   - zcash/librustzcash @ e30517e4335b4d850d35549f8366ba8bf817b17d
%%     (2026-07-10): zcash_keys unstable-frost surface.
%%   - zips PR #895 (open, fetched 2026-07-15).
%%   - C. P. L. Gouvea, C. Komlo, "Re-Randomized FROST", IACR ePrint
%%     2024/436, https://eprint.iacr.org/2024/436 (Secs. 5.1, 7; Remark 1).
%%   - NCC Group public report (2023-10-23); Taurus/ZenGo audit report
%%     (2021-03-23, vendored in the frost and reddsa repos); Least
%%     Authority frost-client/frostd audit announcement (2025-05-01),
%%     https://zfnd.org/frost-demo-audit-frost-client-and-frostd/.
%%   - "The State of FROST for Zcash", zfnd.org, 2025-05-20.
%%   - ZF FROST book: zcash/technical-details.md, zcash/devtool-demo.md,
%%     zcash/server.md (in the frost repo, book/src/).
%%   - Sibling volumes: ironwood-guide.tex (sec:redpallas at line 1308,
%%     sec:redpallas-sec at line 1742), wallet-guide.tex (sec:pczt,
%%     sec:pczt-sign), crosslink-guide.tex (sec:xl-stake-keys).
%%
%% ZIP note: the FROST ZIP is ZIP 312 alone ("FROST for Spend
%% Authorization Multisignatures", Draft, created 2022-08). ZIP 313 is a
%% conventional-fee ZIP (Status: Obsolete, obsoleted by ZIP 317) and is
%% not cited by this volume.
%%
%% Macro note: \ask, \aksym, \nksym, \rk, \F, \Pallas, \Vesta per the
%% ironwood-guide preamble; the volume preamble must define them.
%% Cross-references: sec:fr-keygen and sec:fr-rfc resolve against
%% 02-fr-frost.tex (frost-as-published.tex duplicates both labels;
%% assembly must keep exactly one of the two files); sec:fr-rerand and
%% sec:fr-rerand-suite resolve against 03-fr-rerand.tex.

\section{Deployment and open questions}\label{sec:fr-deploy}

The preceding sections dealt in designs and proofs; this one deals in what
ships. We pin the exact code path a Zcash threshold-custody deployment
exercises today, reconcile it against the draft ZIP text where the two have
diverged, then walk the three custody surfaces this volume targets ---
spend authorisation through PCZT, the ZIP-2005 quantum spending key, and
Crosslink finalizer keys --- and close by sorting every artefact into the
series' three bins (\emph{specified}, \emph{designed-but-unspecified},
\emph{open problem}), each with the concrete trigger on which its
classification must be revisited.

\subsection{The deployed stack}\label{sec:fr-deploy-stack}

Two Zcash Foundation repositories carry the implementation. The
\texttt{frost} workspace (commit \texttt{2016e44ba4}, 2026-04-23, tagged
\texttt{frost-core/v3.0.0}; workspace version 3.0.0, MSRV 1.81) provides
the generic \texttt{frost-core}, the RFC 9591 ciphersuite crates
\texttt{frost-ristretto255}, \texttt{-ed25519}, \texttt{-ed448},
\texttt{-p256}, \texttt{-secp256k1}, and \texttt{-secp256k1-tr}, and the
generic re-randomisation layer \texttt{frost-rerandomized}. Beyond
RFC 9591's two-round protocol the workspace implements trusted-dealer key
generation (the RFC's appendix), the Pedersen-with-knowledge-proof DKG ``as
specified in the original paper'' (ePrint 2020/852,
\S\ref{sec:fr-keygen} --- not in the RFC), Repairable-Threshold-Scheme
share recovery (ePrint 2017/1155), share refresh, and Re-Randomized FROST
(ePrint 2024/436, \S\ref{sec:fr-rerand}). The README declares the crates
``considered stable and feature complete'', with SemVer guarantees.

The \texttt{frost-rerandomized} crate is \texttt{no\_std}; its
\texttt{Randomize} implementations add the randomizer group element to
every verifying share of a \texttt{PublicKeyPackage} and the randomizer
scalar to the signing share of a \texttt{KeyPackage}, whose
\texttt{randomize} documentation warns: ``You MUST NOT reuse the
randomized key package for more than one signing.'' Its README notes that
``the main ciphersuite crates do not re-expose the rerandomization
functions\dots\ The exceptions are the Zcash ciphersuites in
\texttt{reddsa} which do expose the randomized functionality'' --- the
re-randomised layer exists, in shipped form, for Zcash alone.

The \texttt{reddsa} crate (commit \texttt{3792daa95e}, 2026-05-22, version
0.5.2) supplies those ciphersuites: its feature \texttt{frost} $=$
\texttt{["frost-rerandomized", "alloc"]} pulls
\texttt{frost-rerandomized} 3.0.0, and its modules
\texttt{reddsa::frost::redpallas} and \texttt{reddsa::frost::redjubjub}
implement the two ZIP-312 ciphersuites (\texttt{PallasBlake2b512},
\texttt{JubjubBlake2b512}), each with a \texttt{keys} module
(trusted-dealer, DKG, repairable) and a \texttt{rerandomized} submodule
re-exporting \texttt{sign\_with\_randomizer\_seed}, \texttt{aggregate},
\texttt{aggregate\_custom}, \texttt{Randomizer}, and
\texttt{RandomizedParams}. Function \texttt{hash\_randomizer} instantiates
ZIP 312's $H_R$ as $H^{\circledast}$ (the code's \texttt{HStar}) under
the personalisation \texttt{FROST\_RedPallasA} (respectively
\texttt{FROST\_RedJubjubA}), matching \S\ref{sec:fr-rerand-suite}.

The RedPallas module additionally enforces the even-$y$ requirement that
ZIP 2005's amended protocol-spec \S4.2.3 places on $\aksym$
(\S\ref{sec:fr-deploy-qsk}): trait \texttt{keys::EvenY} negates the group
verifying key, every verifying share, the signing shares, and the VSS
commitment coefficients whenever the serialised key has $\tilde{y} = 1$
(the code's evenness test is \texttt{serialized[31] \& 0x80 == 0}), and
the
\texttt{post\_generate} and \texttt{post\_dkg} ciphersuite hooks call
\texttt{into\_even\_y} on fresh key material, so trusted-dealer and DKG
outputs alike land on the even-$y$ representative. Negation commutes with
Shamir sharing because it is linear. The RedJubjub module carries no such
machinery; the constraint is Pallas/Orchard-specific.

\subsection{The randomizer flow: draft ZIP against shipped crate}
\label{sec:fr-deploy-rand}

ZIP 312 (``FROST for Spend Authorization Multisignatures''; owners Gouvea,
Komlo, Connolly; Status Draft; created 2022-08) specifies, at zips commit
\texttt{6961098410}, the randomizer generation
\begin{align*}
\mathtt{signing\_package\_enc} &=
  \mathtt{encode\_signing\_package}(\mathtt{commitment\_list},
  \mathtt{msg}), \\
\mathtt{randomizer} &= H_R\bigl(\mathtt{random\_bytes}(N_s) \,\|\,
  \mathtt{signing\_package\_enc}\bigr),
\end{align*}
with $H_R$ for FROST(Pallas, BLAKE2b-512) being BLAKE2b-512 under the
personalisation \texttt{FROST\_RedPallasA}, output reduced to a Pallas
scalar
(an element of $\F_{p_{\Vesta}}$); $H_2$ is \texttt{Zcash\_RedPallasH},
i.e.\ exactly RedPallas's $H^{\circledast}$
(\S\ref{sec:fr-rerand-suite}). For Jubjub the personalisations are
\texttt{FROST\_RedJubjub\{R,N,M,C,A\}} and \texttt{Zcash\_RedJubjubH},
with base point $G^{\mathsf{Sapling}}$. Round 2 then proceeds as plain
FROST with three substitutions,
\begin{gather*}
sk_i \gets sk_i + \mathtt{randomizer}, \qquad
PK_i \gets PK_i + \mathtt{ScalarBaseMult}(\mathtt{randomizer}), \\
\mathit{group\_public\_key} \gets \mathit{group\_public\_key} +
\mathtt{ScalarBaseMult}(\mathtt{randomizer}),
\end{gather*}
the second applying in share verification. The
Coordinator sends the randomizer over a ``confidential and authenticated
channel'' --- strictly stronger than plain FROST's authenticated-only
channel assumption, because the randomizer links the ceremony to the
on-chain $\rk$.

The deployed crate has moved past this text. Crate
\texttt{frost-rerandomized} 3.0.0 deprecates
\texttt{Randomizer::new(rng, signing\_package)} --- the exact draft-ZIP
flow above, message included --- in favour of
\texttt{Randomizer::new\_from\_commitments}: the Coordinator draws an
$N_s$-byte \texttt{randomizer\_seed}, sends the \emph{seed}, and each
participant regenerates
\[
\mathtt{randomizer} \;=\; \mathtt{hash\_randomizer}\bigl(\mathtt{seed}
\,\|\, \mathtt{encode\_group\_commitments}(\mathtt{commitments})\bigr)
\]
locally via \texttt{regenerate\_from\_seed\_and\_commitments}. The message
is no longer an input, and --- the crate's stated rationale ---
participants no longer need to fully trust the Coordinator's RNG, since
their own round-one commitments enter the hash. In the paper's terms
(ePrint 2024/436, Definition 5 and Remark 1), the commitment encoding
plays the role of the optional contributory string $\mathrm{aux}$, which
models ``the transcript of some external protocol''; the paper's own
\S7 implementation notes describe the older design, with the message
still included, in the same terms.

The gap between draft and crate is being closed in zips PR \#895
(``[ZIP 312] Specify key generation, change randomizer handling''; opened
August 2024, 21 commits, still open as of 2026-07-15, awaiting review by
daira and nuttycom). It adds a key-generation specification, revises the
randomizer handling (``the previous approach couldn't work using messages
as input''), integrates the ZIP-2005 $\mathsf{qsk}$ material, renames the
scheme consistently to ``Re-Randomized FROST'', and adds Daira-Emma
Hopwood as a ZIP owner. Its commit history also retires the draft's
serialisation reference --- the ZF FROST book's moving-target
serialisation-format page --- in favour of RFC 9591's serialisation
function. Until it merges, this volume presents the crate flow ---
message omission included --- as the observed fact of the pinned
commit, draft lag rather than a conformance defect, and the ZIP text as
in flight.
%% TODO: when #895 merges, re-derive this subsection from the merged
%% text; it is the single sharpest revisit trigger for the chapter.

The trust geometry is worth stating exactly. The paper's \S5.1 holds that
the central randomizer-choosing party is ``trusted with preserving privacy
of the scheme'' but ``not trusted with security'', and that removing it
would require a DKG-style contributory sub-protocol at extra rounds and
complexity. ZIP 312's threat model matches: the Coordinator is trusted
with transaction privacy and unlinkability but cannot forge without
$\mathtt{MIN\_PARTICIPANTS}$ approvals, and every key-share holder is
trusted with transaction privacy. Non-requirements are equally explicit:
ZIP 312 does not remove the Coordinator role, does not (at this commit)
provide key generation, and places network privacy out of scope.

\subsection{PCZT as the integration surface}\label{sec:fr-deploy-pczt}

ZIP 374's Motivation names ``multiparty transactions --- threshold signing
(for example FROST [RFC 9591])'' among the problems the PCZT format
solves; the \emph{Wallet Guide}, \S``Partially created transactions
(PCZT)'', documents the format and its role graph. The insertion point for
FROST is the Signer role (\emph{Wallet Guide}, \S``Signing''): to sign an
Orchard spend the action's re-randomisation scalar $\alpha$ must be
present, the local-key path computes $\rsk := \ask + \alpha$ ---
the re-randomisation of the \emph{Ironwood Guide}, \S``Authorisation of
the spend: RedPallas'' --- and checks
$\mathsf{VerificationKey}(\rsk) = \rk$, and the
\emph{external} path, \texttt{apply\_orchard\_signature}, accepts a
signature if and only if it verifies under the action's $\rk$ over the
shielded sighash (\texttt{InvalidExternalSignature} otherwise). A
FROST-aggregated RedPallas signature is exactly such an external
signature; for it to verify, the Coordinator must use the PCZT action's
$\alpha$ as the randomizer, so that the re-randomised group key
$\mathsf{RandPK}$ produces equals that action's $\rk$.

The end-to-end flow has been demonstrated (ZF FROST book,
\texttt{zcash/devtool-demo.md}):
\begin{enumerate}
\item \texttt{zcash-devtool pczt create} produces the PCZT;
\item the \texttt{zcash-sign} tool (from
  ZcashFoundation/\allowbreak frost-zcash-demo, now \texttt{frost-tools})
  prints a SIGHASH and a \texttt{Randomizer} from the transaction plan;
\item \texttt{frost-client} runs the FROST protocol over \texttt{frostd}
  with that SIGHASH as the message and that randomizer;
\item \texttt{zcash-sign sign} writes the aggregated signature into
  \texttt{frost\_pczt.signed};
\item \texttt{zcash-devtool pczt combine} merges the signed and proven
  PCZTs;
\item \texttt{zcash-devtool pczt send} broadcasts.
\end{enumerate}
Here \texttt{zcash-devtool} builds on \texttt{zcash\_keys} with the
features \texttt{"orchard"}, \texttt{"unstable"}, and
\texttt{"unstable-frost"}. The \texttt{unstable-frost} feature ($=$
\texttt{["orchard"]}) exposes the constructor
\texttt{UnifiedFullViewingKey::}\allowbreak\texttt{from\_orchard\_fvk}
--- a UFVK from an Orchard FVK alone, since a FROST wallet has no
unified spending key --- and the \texttt{zcash\_keys} changelog (0.3.0,
2024-08-19) frames the flag as existing ``to temporarily expose API
features that are needed
specifically when creating FROST threshold signatures'', to be removed
``once key derivation for FROST has been fully specified and
implemented''. Repository \texttt{librustzcash} itself, that is,
classifies FROST key derivation as not fully specified.

One check remains outside every specification: ZIP 312 requires that
signers ``MUST check that the given SIGHASH matches the data sent from the
Coordinator, or compute the SIGHASH themselves'', and explicitly leaves
the mechanism out of scope. In the demonstrated flow the SIGHASH travels
from \texttt{zcash-sign} to the participants through
\texttt{frost-client}; by what mechanism a share-holding participant
gains independent confidence in it is unspecified.
%% TODO: PCZT v2 (v6 transactions, ZIP 230/229 era) may change the
%% Signer/alpha surface the Coordinator consumes; recheck this
%% subsection against the Wallet Guide's sec:pczt-v2 material when v2
%% ships.

\subsection{Wallet-integration constraints}\label{sec:fr-deploy-wallet}

The ZF FROST book's technical-details chapter records the constraints a
FROST-capable Zcash wallet inherits. FROST cannot cover transparent
inputs, which authorise by ECDSA; the book accordingly suggests supporting
Orchard only (``Orchard is simpler to handle''). With a DKG-generated
$\aksym$, the remaining Orchard key components derive from $\aksym$ plus
either an $\mathsf{sk}$ used only for $\nksym$ and $\mathsf{rivk}$, or
independently generated $\nksym$ and $\mathsf{rivk}$. A seed phrase alone
can never restore a FROST wallet: restoration needs the key share, all
verifying shares, the participant identifiers, and the FVK components, so
share backup is a JSON-file affair and share recovery falls to the
Repairable Threshold Scheme (\S\ref{sec:fr-deploy-stack}).

The Foundation's own status report (``The State of FROST for Zcash'',
2025-05-20) names wallet integration the missing piece (``The missing
piece is for wallets to integrate FROST''): no
mainstream wallet ships FROST, and ``technical-inclined users can use
FROST with Zcash today using our \texttt{frost-client} command line
tool''. The same post lists the ``lack of a standardized key generation
specification for FROST'' as the interoperability limiter, flags metadata
protection (e.g.\ Tor) as needing work, and offers ZF as a candidate
operator of a production \texttt{frostd}.

\subsection{Transport: frostd and frost-client}
\label{sec:fr-deploy-transport}

The reference transport is \texttt{frostd}, a JSON-HTTP session server:
clients authenticate with key pairs, the Coordinator creates sessions
naming the participants' public keys, and all FROST messages are
end-to-end encrypted, so the server itself is untrusted. The CLI
\texttt{frost-client} drives it; both live in
ZcashFoundation/\allowbreak frost-tools (renamed from
\texttt{frost-zcash-demo}). Least Authority, as Zcash Ecosystem Security
Lead, audited \texttt{frost-client} and \texttt{frostd} (announcement
2025-05-01) with no
high-severity findings, all findings addressed on \texttt{main}. No ZIP
specifies this transport; ZIP 312 places network privacy out of scope
(\S\ref{sec:fr-deploy-rand}).

\subsection{Audit coverage, precisely}\label{sec:fr-deploy-audit}

Two audits exist, and neither covers the code path Zcash custody would
run. NCC Group audited the v0.6.0 release (commit \texttt{5fa17ed}) of
\texttt{frost-core}, \texttt{frost-ed25519}, \texttt{frost-ed448},
\texttt{frost-p256}, \texttt{frost-secp256k1}, and
\texttt{frost-ristretto255} --- including trusted-dealer and DKG key
generation and FROST signing --- with all findings addressed and reviewed
by NCC (public report 2023-10-23). The repository README is explicit:
``This does not include \texttt{frost-secp256k1-tr} and rerandomized
FROST.'' Earlier, the 2021 ``Audit of FROST implementation'' (JP Aumasson
and Adrien Hamelink, Taurus; Omer Shlomovits, ZenGo; dated 2021-03-23)
covered the pre-rewrite implementation in the RedJubjub repository's
\texttt{src/frost.rs} ($\approx 404$ lines at commit \texttt{2ebc08f},
2021-02-25) --- two-round FROST with a trusted dealer only, RedJubjub
only --- concluding ``We did not find any major security issue.''

The composition, therefore: \texttt{frost-rerandomized} plus the
\texttt{reddsa} RedPallas/RedJubjub ciphersuites --- the exact stack of
\S\ref{sec:fr-deploy-stack} --- has no completed third-party audit, and
the only audit ever to touch a Zcash FROST ciphersuite predates the
\texttt{frost-core} rewrite.
%% Audit landscape rechecked 2026-07-15: frost README @2016e44ba4 still
%% scopes the NCC audit to v0.6.0 and excludes rerandomized FROST; the
%% Least Authority announcement (2025-05-01) covers the frost-client and
%% frostd tools only; no audit touches frost-rerandomized or the reddsa
%% suites. A print-time recheck is mandated by the revisit trigger in
%% sec:fr-deploy-class.

\subsection{Custody of the ZIP-2005 quantum spending key}
\label{sec:fr-deploy-qsk}

ZIP 2005 (``Ironwood Quantum Recoverability''; owners Daira-Emma Hopwood
and Jack Grigg; Status Proposed, Category Consensus; activation at the
NU6.3 activation height) is the second custody surface. Its ``Usage with
FROST'' section derives $\aksym$ jointly, from the participants' shares of
$\ask$, via FROST DKG, and then constrains the ceremony: participants
``MUST privately agree on a value $\mathsf{sk}$, and then use it with
$\mathsf{use\_qsk} = \mathsf{true}$ to derive $\nksym$,
$\mathsf{qsk}$, $\mathsf{qk}$, and (using the $\aksym$ output by the
DKG protocol) $\mathsf{rivk\_ext}$''; the protocol MUST ensure all
participants obtain the same $\aksym$, $\nksym$, and
$\mathsf{rivk\_ext}$. The derivations, verbatim from the ZIP:
\begin{align*}
\mathsf{H}^{\mathsf{qsk}}(\mathsf{sk}) &=
  \mathrm{truncate}_{32}\bigl(\mathsf{PRFexpand}_{\mathsf{sk}}
  ([\mathtt{0x0C}])\bigr), \\
\mathsf{H}^{\mathsf{qk}}(\mathsf{qsk}) &=
  \mathrm{BLAKE3.derive\_key}(\texttt{"Zcash ZIP 2005 qk-derivation
  v1"}, \mathsf{qsk}, 32), \\
\mathsf{H}^{\mathsf{rivk\_ext}}_{\mathsf{qk}}(\aksym, \nksym) &=
  \mathrm{ToScalar}^{\mathsf{Orchard}}\bigl(
  \mathsf{PRFexpand}_{\mathsf{qk}}([\mathtt{0x0D}] \,\|\,
  \mathrm{I2LEOSP}_{256}(\aksym) \,\|\,
  \mathrm{I2LEOSP}_{256}(\nksym))\bigr),
\end{align*}
with the $\mathsf{PRFexpand}$ domain-separator bytes
$\mathtt{0x0A}$--$\mathtt{0x0D}$ registered by ZIP 2005. In the
$\mathsf{use\_qsk} = \mathsf{true}$ branch of the amended protocol-spec
\S4.2.3, the wallet obtains $\aksym^{\mathbb{P}}$ ``by any suitably
secure method that ensures the last bit of
$\mathrm{repr}_{\mathbb{P}}(\aksym^{\mathbb{P}})$ is 0'' --- the
requirement the \texttt{reddsa} \texttt{EvenY} machinery of
\S\ref{sec:fr-deploy-stack} enforces --- and sets $\aksym =
\mathrm{Extract}_{\mathbb{P}}(\aksym^{\mathbb{P}})$, $\mathsf{qsk} =
\mathsf{H}^{\mathsf{qsk}}(\mathsf{sk})$, $\mathsf{qk} =
\mathsf{H}^{\mathsf{qk}}(\mathsf{qsk})$, and $\mathsf{rivk} =
\mathsf{H}^{\mathsf{rivk\_ext}}_{\mathsf{qk}}(\aksym, \nksym)$.
The ZIP's list of admissible generation methods for
$\aksym^{\mathbb{P}}$ explicitly includes ``FROST distributed key
generation [zip-0312-key-generation] in which the corresponding
spend-authorization material is $t$-of-$n$ shared among the
participants''.

The threat model shifts accordingly. The host wallet in a multi-signer
FROST setup is the ZIP-312 Coordinator; it holds $\nksym$, $\aksym$, and
$\mathsf{qk}$. With FROST and hardware wallets, spend authorisation breaks
only by (1) possession of $\mathsf{qsk}$ \emph{plus} finding discrete
logarithms on Pallas --- and ``the adversary is an authorized FROST
signer --- they hold a copy of $\mathsf{qsk}$ by construction'' is the
ZIP's variant 1.a --- or (2) breaking RedDSA, FROST, the DKG, or the
proving system. Each participant MUST treat $\mathsf{qsk}$ as a secret
held at least as securely and reliably as $\ask$: it is not needed until
the Recovery Protocol, but losing it can lose funds once the current
Orchard protocol is disabled. A quantum adversary holding only
$\mathsf{qsk}$ --- which every participant holds --- may steal funds, so
the ZIP RECOMMENDS that once a fully post-quantum threshold
multisignature protocol exists, all FROST-held funds migrate to its pool.
FROST's key advantage --- signing with shares without ever reconstructing
$\ask$ --- is retained through the Recovery Protocol, which continues to
require a RedDSA signature.

One tension is unresolved. FROST DKG exists precisely to avoid a
commonly-known secret, yet the ceremony above requires the participants to
``privately agree on a value $\mathsf{sk}$'' from which $\nksym$ and
$\mathsf{qsk}$ derive --- and neither ZIP 2005 nor the crates give a
protocol for that agreement. ZIP 2005's Deployment section points at
PR \#895 (``There is no significant existing deployment of FROST, so we
can write this into ZIP 312 from the start''), so the resolution is
expected there: as of 2026-07-15 the PR's commit history drafts a
contributory $\mathsf{sk}$-generation section (2026-03-05), still
unmerged, so the shared-$\mathsf{sk}$ agreement stays \emph{designed
but unspecified} with \#895's merge as its trigger.

\subsection{Crosslink finalizer keys}\label{sec:fr-deploy-crosslink}

The third surface is prospective only. The \emph{Crosslink Guide},
\S``Finalizer keys'', records that the prototype uses single ed25519 keys
per finalizer (ed25519-zebra), with 76-byte votes signed by appending a
64-byte ed25519 signature over those 76 bytes, and that no FROST or
threshold-signing design for finalizer
keys is sourced anywhere in the Crosslink corpus (negative search over
tfl-book, crosslink-protocol-guide, crosslink-spec, and zebra-crosslink,
dated 2026-07-15 in that guide); key rotation, revocation, threshold
custody, and initial-roster formation are classified open problems in
every design source.

We record one observation of our own, clearly marked as such: finalizer
votes are plain ed25519, with no key re-randomisation, so the fitting
ciphersuite would be RFC 9591's FROST(Ed25519, SHA-512) via the
NCC-audited
\texttt{frost-ed25519} crate, with no re-randomisation layer --- finalizer
keys are deliberately long-lived roster identities, not unlinkable
one-time keys, so the entire apparatus of \S\ref{sec:fr-rerand} is
unnecessary there. No published design says this; it is this volume's
inference from the vote format and the audit table.
%% TODO: revisit when Shielded Labs publishes any finalizer
%% key-management or rotation design; none exists as of 2026-07-15.

\subsection{Classification and revisit triggers}
\label{sec:fr-deploy-class}

Per the series' three-way scheme:

\begin{itemize}
\item \emph{Specified.} FROST as in RFC 9591 (\S\ref{sec:fr-rfc});
  Re-Randomized FROST for spend-authorisation custody, per ZIP 312 at
  Draft status --- PR \#895's in-flight randomizer and key-generation
  changes are its named revisit trigger and do not move the bin
  (\S\ref{sec:fr-deploy-rand}); the
  \texttt{frost-core}/\texttt{frost-rerandomized}/\texttt{reddsa} API
  surface at the pinned versions (\S\ref{sec:fr-deploy-stack}); the PCZT
  Signer external-signature path (\S\ref{sec:fr-deploy-pczt}).
\item \emph{Designed-but-unspecified.} FROST key generation for Zcash
  (out of scope at the pinned ZIP-312 commit; the crates implement the
  PedPoP DKG of ePrint 2020/852, but no ZIP specifies it); the ZIP-2005
  $\mathsf{use\_qsk}$ FROST
  key-generation constraints (ZIP Proposed, referencing the unmerged
  \#895; the $\mathsf{sk}$-agreement step has no protocol,
  \S\ref{sec:fr-deploy-qsk}); the transport
  (\texttt{frostd} exists and is audited, but no ZIP covers it,
  \S\ref{sec:fr-deploy-transport}).
\item \emph{Open problem.} Crosslink finalizer threshold custody
  (\S\ref{sec:fr-deploy-crosslink}); post-quantum threshold signing for
  the $\mathsf{qsk}$-era migration (\S\ref{sec:fr-deploy-qsk});
  hardware-wallet FROST share custody --- ZIP 2005's threat model assumes
  hardware-wallet FROST signers, but no device implements a share signer.
\end{itemize}

The classifications above expire on identifiable events. In rough order
of expected impact:

\begin{itemize}
\item \emph{PR \#895 merges.} The randomizer flow, the key-generation
  specification, the $\mathsf{sk}$-agreement protocol, and the
  ZIP-312/ZIP-2005 integration all re-derive from the merged text; the
  serialisation reference it pins replaces the ZF-FROST-book moving
  target (\S\ref{sec:fr-deploy-rand}).
\item \emph{An audit of the re-randomised stack is announced.} The
  ``no completed third-party audit'' claim of
  \S\ref{sec:fr-deploy-audit} must be rechecked immediately before any
  printing in any case.
\item \emph{NU6.3 approaches.} Whether $\mathsf{use\_qsk} =
  \mathsf{true}$ key generation becomes normative for all new FROST
  wallets, and whether \texttt{zcash\_keys}' \texttt{unstable-frost}
  surface is replaced by a specified derivation, both resolve with the
  upgrade timeline (\S\ref{sec:fr-deploy-pczt},
  \S\ref{sec:fr-deploy-qsk}).
\item \emph{PCZT v2 ships.} The Signer/$\alpha$ surface the Coordinator
  consumes may change with v6 transactions
  (\S\ref{sec:fr-deploy-pczt}).
\item \emph{A finalizer key-management design is published.} The
  Crosslink bin moves, and the \texttt{frost-ed25519} observation is
  either confirmed or discarded (\S\ref{sec:fr-deploy-crosslink}).
\item \emph{A post-quantum threshold multisignature protocol exists.}
  ZIP 2005's RECOMMENDED migration of FROST-held funds becomes
  actionable (\S\ref{sec:fr-deploy-qsk}).
\end{itemize}

\end{document}
